% =====================================================================================
% ch03_prelim.tex —— 第 3 章　因子化的数学与场论工具（writer W03）
% 约定：WRITER_SPEC_v2 §3；公式来源 VERIFIED_FORMULAS_v2（VF-…，见行内注释）；
% 独立核算脚本 writing/work/W03/s1–s6（日志 writing/work/W03/logs/）。
% =====================================================================================
\chapter{因子化的数学与场论工具}\label{ch:prelim}

因子化定理断言：高能过程的截面可以写成只依赖于硬标度 $Q$ 的短程系数与只依赖于强子结构的长程矩阵元的卷积，误差为 $Q$ 的负幂次。要把这句话变成可以证明的数学陈述，先要回答几个看似无关的问题：高速运动的强子的哪些动量分量大、哪些小？怎样把不同时空点的夸克场组合成规范不变的算符，高能夸克在软胶子"看来"又是什么？维数正规化中的 $1/\varepsilon$ 为什么有时来自紫外、有时来自红外？单圈计算中的 $(1-z)^{-1-2\varepsilon}$ 应当怎样理解？矩能确定多少 $x$ 空间的信息？截面与前向振幅的虚部是什么关系？本章依次回答这些问题，并与\cref{sec:qcd-quantization}的约定一起固定全书的记号。

除一处以外，本章的陈述或者是纯数学（不带标签），或者是由定义直接推出的严格结论（\tagR，$R_{\mathrm{def}}$），或者是关于 Feynman 图的严格陈述（\tagR，$R_{\mathrm{PT}}$）；例外是\cref{rem:prelim-cutting-dr}：无质量维数正规化中含鬼场的割图规则只有全阶论证（\tagA）。涉及强子态的只有光学定理（\cref{prop:prelim-optical}），其强子版本需要\cref{asm:intro-wightman}。较长的证明在附录~A、B、C、I、J、K。

\begin{tcolorbox}[enhanced,breakable,title=本章证明了什么／没有证明什么,fonttitle=\bfseries,
  colback=gray!3!white,colframe=gray!55!black,boxrule=0.6pt,arc=2pt,
  left=5pt,right=5pt,top=3pt,bottom=3pt,before skip=6pt,after skip=6pt]
\textbf{证明了：}\cref{lem:prelim-boost}（\tagR，$R_{\mathrm{def}}$）；\cref{prop:prelim-wilson-covariance}、\cref{prop:prelim-wilson-rules}（\tagR，$R_{\mathrm{def}}$；后者的完整证明见附录~B）；\cref{lem:prelim-eikonal-identity}、\cref{lem:prelim-scaleless}、\cref{lem:prelim-plus-expansion}、\cref{lem:prelim-mellin}（数学；后两者的证明见附录~A）；\cref{thm:prelim-cutting-rules}（\tagR，$R_{\mathrm{PT}}$，对有质量或已正规化的图）；\cref{prop:prelim-unitarity}（\tagR，$R_{\mathrm{PT}}$）；\cref{prop:prelim-optical}（\tagR，$R_{\mathrm{def}}$；强子态需\cref{asm:intro-wightman}）。\par
\textbf{给出全阶论证：}\cref{rem:prelim-cutting-dr}（\tagA：无质量维数正规化中含鬼场的割图规则）。\par
\textbf{没有证明：}带插入的光锥割线位置恒等式与闭时路形式的幺正性（附录~O 的\cref{lem:appO-cut-position}、\cref{lem:appO-unitarity}）；单个割图并不满足幺正性，只有对全部割线求和才行（\cref{sec:prelim-cuts}末的常见误解）。
\end{tcolorbox}

% -------------------------------------------------------------------------------------
\section{光锥坐标、Sudakov 分解与 Lorentz 推动}\label{sec:prelim-lightfront}

质量为 $M$、能量 $E\gg M$ 的强子沿 $+z$ 运动时，$E+P^3\approx2E$，而 $E-P^3=M^2/(E+P^3)\approx M^2/(2E)$，二者相差 $(2E/M)^2$ 倍；与它几乎平行的部分子也有同样的层级。因此天然的坐标是光锥坐标 $v^\pm=(v^0\pm v^3)/\sqrt2$：沿 $z$ 轴的推动只是把 $v^\pm$ 乘以 $\ee^{\pm\eta}$（\cref{lem:prelim-boost}）。后文的幂次计数正建立在这一事实之上。

\begin{definition}[光锥坐标与 Sudakov 分解]\label{def:prelim-lightfront}
(a) 对四矢量 $v^\mu=(v^0,v^1,v^2,v^3)$ 定义光锥分量 $v^\pm=(v^0\pm v^3)/\sqrt2$ 与横向二维矢量 $\mathbf v_{\mathrm T}=(v^1,v^2)$，记 $v=(v^+,v^-,\mathbf v_{\mathrm T})$。在度规 $(+,-,-,-)$ 下
\begin{equation}\label{eq:prelim-lc-dot}
v\cdot w=v^+w^-+v^-w^+-\mathbf v_{\mathrm T}\cdot\mathbf w_{\mathrm T},\qquad v^2=2v^+v^--\mathbf v_{\mathrm T}^2 ,
\end{equation}
即 $g^{+-}=g^{-+}=1$，$g^{++}=g^{--}=0$，$v_\pm=v^\mp$；积分测度为 $\dd^4k=\dd k^+\dd k^-\dd^2\mathbf k_{\mathrm T}$（$(k^0,k^3)\mapsto(k^+,k^-)$ 的 Jacobi 行列式的绝对值为 $1$）。% VF-CONV-1
\par
(b) 类光矢量 $n_A=(1^+,0^-,\mathbf 0_{\mathrm T})$ 与 $n_B=(0^+,1^-,\mathbf 0_{\mathrm T})$ 满足
\begin{equation}\label{eq:prelim-nAnB}
n_A\cdot v=v^-,\qquad n_B\cdot v=v^+,\qquad n_A^2=n_B^2=0,\qquad n_A\cdot n_B=1 .
\end{equation}
强子 $A$ 沿 $+z$ 运动，$P_A=(P_A^+,M_A^2/(2P_A^+),\mathbf 0_{\mathrm T})$；强子 $B$ 沿 $-z$ 运动，$P_B=(M_B^2/(2P_B^-),P_B^-,\mathbf 0_{\mathrm T})$。
\par
(c) Sudakov 分解：取类光矢量 $p_A=p_A^+n_A$，$p_B=p_B^-n_B$，任一四矢量唯一地写成
\begin{equation}\label{eq:prelim-sudakov}
k^\mu=\alpha\,p_A^\mu+\beta\,p_B^\mu+k_{\mathrm T}^\mu,\qquad \alpha=\frac{k^+}{p_A^+},\quad \beta=\frac{k^-}{p_B^-},\quad k_{\mathrm T}^\mu=(0,\mathbf k_{\mathrm T},0),
\end{equation}
于是 $k^2=2\alpha\beta\,p_A\cdot p_B-\mathbf k_{\mathrm T}^2$，$p_A\cdot p_B=p_A^+p_B^-$。
\par
(d) 快度 $y=\frac12\ln(k^+/k^-)$（$k^\pm>0$）。质量为 $m$ 的在壳粒子有 $k^\pm=m_{\mathrm T}\ee^{\pm y}/\sqrt2$，$m_{\mathrm T}=(m^2+\mathbf k_{\mathrm T}^2)^{1/2}$。
% uses: none
\end{definition}

$n_B$ 即 Collins~\cite{Collins2011}的 $w_2$；第~15~章 SCET 的 $n=(1,0,0,1)$、$\bar n=(1,0,0,-1)$ 与 $n_A,n_B$ 的换算见\cref{sec:appA-lightfront}。强子 $A$ 的部分子分布中的 Wilson 线沿 $n_B$，只含 $n_B\cdot A=A^+$（\cref{def:pdf-quark}，第~8~章）；为什么是 $A^+$，由下面的推动引理回答。

\begin{lemma}[推动下的标度 \tagR]\label{lem:prelim-boost}
（$R_{\mathrm{def}}$：由算符定义直接推出。）设 $\Lambda(\eta)$ 是沿 $+z$ 方向、快度为 $\eta$ 的推动：$v'^0=v^0\cosh\eta+v^3\sinh\eta$，$v'^3=v^0\sinh\eta+v^3\cosh\eta$，$\mathbf v'_{\mathrm T}=\mathbf v_{\mathrm T}$。
\par
(a) $v'^\pm=\ee^{\pm\eta}v^\pm$，$\mathbf v'_{\mathrm T}=\mathbf v_{\mathrm T}$，快度平移 $y\to y+\eta$；矢量场的分量 $A^\pm$、$\mathbf A_{\mathrm T}$ 按同样方式变换。
\par
(b) 令 $\gamma^\pm=(\gamma^0\pm\gamma^3)/\sqrt2$，$P_+=\frac12\gamma^-\gamma^+$，$P_-=\frac12\gamma^+\gamma^-$。它们是互补的正交投影算符，$P_\pm=\frac12(1\pm\gamma^0\gamma^3)$。在旋量表示中实现 $\Lambda(\eta)$ 的矩阵（$S^{-1}\gamma^\mu S=\Lambda^\mu{}_\nu\gamma^\nu$）为
\begin{equation}\label{eq:prelim-boost-spinor}
S(\eta)=\exp\Big(\frac\eta2\gamma^0\gamma^3\Big)=\ee^{\eta/2}P_++\ee^{-\eta/2}P_- .
\end{equation}
因此推动后的 Dirac 场 $\psi'(x)=S(\eta)\psi(\Lambda^{-1}x)$ 的"好分量" $P_+\psi$ 与"坏分量" $P_-\psi$ 分别乘以 $\ee^{\eta/2}$ 与 $\ee^{-\eta/2}$：$P_\pm\psi'(\Lambda x)=\ee^{\pm\eta/2}P_\pm\psi(x)$。
\uses{def:prelim-lightfront}
\end{lemma}

\begin{proof}
完整计算见\cref{sec:appA-lightfront}与\cref{sec:appA-dirac}；要点如下。(a) 由 $\cosh\eta\pm\sinh\eta=\ee^{\pm\eta}$ 直接得到。(b) $(\gamma^0\gamma^3)^2=-(\gamma^0)^2(\gamma^3)^2=1$，故 $S(\eta)=\cosh\frac\eta2+\gamma^0\gamma^3\sinh\frac\eta2=\ee^{\eta/2}P_++\ee^{-\eta/2}P_-$；附录~A 证明 $\gamma^+=P_-\gamma^+P_+$、$\gamma^-=P_+\gamma^-P_-$ 且 $\gamma^{1,2}$ 与 $P_\pm$ 对易，由此 $S^{-1}\gamma^\pm S=\ee^{\pm\eta}\gamma^\pm$，$S^{-1}\gamma^{1,2}S=\gamma^{1,2}$；最后 $P_\pm S=\ee^{\pm\eta/2}P_\pm$。
\end{proof}

\begin{example}[共线动量与好/坏分量的标度]\label{ex:prelim-boost-scaling}
静止的强子中部分子动量的各光锥分量都是 $m$ 的量级：$k_{\mathrm{rest}}=m\,(a,b,\mathbf c)$，$a,b,|\mathbf c|=O(1)$。推动到 $P^+\sim Q$ 所需的快度满足 $\ee^\eta=Q/m\equiv1/\lambda$，由\cref{lem:prelim-boost}(a)，
\begin{equation}\label{eq:prelim-collinear-scaling}
k=\big(\ee^{\eta}ma,\ \ee^{-\eta}mb,\ m\mathbf c\big)=Q\,\big(a,\ \lambda^2b,\ \lambda\mathbf c\big),
\end{equation}
正是 $A$-共线标度 $Q(1,\lambda^2,\lambda)$。取 $m=1\,\mathrm{GeV}$、$Q=100\,\mathrm{GeV}$，则 $\lambda=0.01$，$\eta=\ln100\approx4.61$，$(1,1,1)\,\mathrm{GeV}$ 变为 $(100,\,0.01,\,1)\,\mathrm{GeV}$。由\cref{eq:prelim-boost-spinor}，好分量放大 $\lambda^{-1/2}=10$ 倍，坏分量缩小为 $1/10$，二者之比增大 $100$ 倍；沿 $+z$ 运动的无质量夸克只有好分量，因为 $\sla p=p^+\gamma^-$ 与 $\sla pu=0$ 给出 $\gamma^-u=0$，即 $P_-u=0$。胶子场中 $A^+$ 放大 $100$ 倍，$A^-$ 缩小为 $1/100$，$\mathbf A_{\mathrm T}$ 不变：快速强子的胶子场以 $A^+=n_B\cdot A$ 为主。
\uses{lem:prelim-boost}
\end{example}

由此，部分子分布中只出现 $A^+$ 的 Wilson 线；在幂次计数中，与共线子图相连、沿 $A^+$ 极化的纵向（K）胶子不带幂次压低，可以任意多地出现在领头幂次中，须由 Ward 恒等式求和成 Wilson 线（\cref{def:ir-scalings}、\cref{prop:ir-power-counting}，第~4~章）。\cref{fig:prelim-lightcone}画出本节的两个事实：推动沿不变双曲线移动点；在对数光锥坐标下，推动是沿反对角线的平移。

%%% BEGIN fig:prelim-lightcone
\begin{figure}[htbp]
\centering
\begin{tikzpicture}[fd,
  fa lc/.style={draw=gray!60!black,line width=0.6pt,dash pattern=on 3pt off 2pt},
  fa hyp/.style={draw=black,line width=0.8pt,dash pattern=on 0.8pt off 1.6pt},
  fa grid/.style={draw=gray!60,line width=0.5pt,dash pattern=on 0.8pt off 1.6pt},
  fa pt/.style={circle,fill=black,draw=none,inner sep=0pt,minimum size=4.2pt},
  fa boost/.style={draw=red!60!black,line width=0.9pt,-{Stealth[length=5pt,width=4.5pt]}},
  fa leg/.style={draw=gray!50!black,line width=0.45pt},
  fa lab/.style={font=\footnotesize},
  fa wlab/.style={font=\footnotesize,fill=white,inner sep=1pt},
  fa cap/.style={font=\footnotesize}]
  % ===================== (a) the (z,t) plane =====================
  \begin{scope}
    % world-band of a hadron moving fast along +z (slope dt/dz = 1/v, v = 0.83)
    \draw[hadron blob,line width=0.5pt] (1.35,-0.45) -- (1.82,-0.45) -- (3.99,2.15) -- (3.52,2.15) -- cycle;
    \draw[fa leg] (1.9,-0.36) -- (2.35,-0.62);
    \node[fa lab,anchor=north west,align=left,inner sep=1pt] at (2.3,-0.55) {沿 $+z$ 快速运动\\的强子世界带};
    % axes and light-cone axes x^pm = (t pm z)/sqrt2
    \draw[lcaxis] (-2.6,0) -- (4.45,0) node[fa lab,below] {$z$};
    \draw[lcaxis] (0,-0.7) -- (0,3.0) node[fa lab,left] {$t$};
    \draw[fa lc,-{Stealth[length=4pt,width=3.4pt]}] (-0.6,-0.6) -- (2.75,2.75) node[fa lab,anchor=south west,inner sep=1pt] {$x^+$};
    \draw[fa lc,-{Stealth[length=4pt,width=3.4pt]}] (0.6,-0.6) -- (-2.75,2.75) node[fa lab,anchor=south east,inner sep=1pt] {$x^-$};
    % invariant hyperbola t^2 - z^2 = c and a boost of rapidity eta along it
    \draw[fa hyp] plot[domain=-1.2:1.2,samples=60,variable=\t] ({1.1*sinh(\t)},{1.1*cosh(\t)});
    \node[fa lab,anchor=south east,inner sep=1pt] at (-0.2,2.12) {$t^2-z^2=c$};
    \draw[fa leg] (-0.62,2.12) -- (-0.8,1.4);
    \draw[fa boost] plot[domain=0.06:0.74,samples=20,variable=\t] ({1.1*sinh(\t)},{1.1*cosh(\t)});
    \node[fa pt] at (0,1.1) {};
    \node[fa pt] at (0.977,1.471) {};
    \node[fa lab,text=red!55!black,anchor=north west] at (0.55,1.12) {推动 $\eta$};
    \node[fa cap] at (0.9,-1.85) {(a)\ $(z,t)$ 平面};
  \end{scope}
  % ===================== (b) the (ln v^+, ln v^-) plane =====================
  \begin{scope}[xshift=6.9cm]
    \begin{scope}
      \clip (0.05,0.05) rectangle (4.3,3.45);
      % lines of constant v^+ v^-  (anti-diagonals, slope -1)
      \foreach \c in {1.4,2.6,3.8,5.0,6.2} \draw[fa lc] (\c-3.6,3.6) -- (\c+0.2,-0.2);
      % lines of constant rapidity y = (1/2) ln(v^+/v^-)  (diagonals, slope +1)
      \foreach \d in {-2.4,-1.2,0,1.2,2.4} \draw[fa grid] (-0.2,-0.2+\d) -- (4.5,4.5+\d);
    \end{scope}
    \draw[lcaxis] (0,0) -- (4.55,0) node[fa lab,below left=1pt and -2pt] {$\ln v^+$};
    \draw[lcaxis] (0,0) -- (0,3.7) node[fa lab,left] {$\ln v^-$};
    % the boost: +eta in ln v^+, -eta in ln v^-
    \draw[fa leg] (1.5,2.3) -- (2.6,2.3) -- (2.6,1.2);
    \node[fa wlab,anchor=south] at (2.05,2.33) {$+\eta$};
    \node[fa wlab,anchor=west] at (2.64,1.75) {$-\eta$};
    \node[fa pt] (P) at (1.5,2.3) {};
    \node[fa pt] (Q) at (2.6,1.2) {};
    \draw[fa boost,shorten >=2pt,shorten <=2pt] (P) -- (Q);
    \node[fa wlab,text=red!55!black,anchor=north east] at (1.95,1.68) {推动 $\eta$};
    \node[fa wlab,anchor=south east] at (1.44,2.36) {$(\ln v^+,\ln v^-)$};
    \node[fa wlab,anchor=north] at (2.6,1.1) {$(\ln v^++\eta,\ \ln v^--\eta)$};
    % legend for the two line families
    \draw[fa lc] (0.1,-0.75) -- (0.75,-0.75);
    \node[fa lab,anchor=west] at (0.8,-0.75) {等 $v^+v^-$（反对角线）};
    \draw[fa grid] (0.1,-1.15) -- (0.75,-1.15);
    \node[fa lab,anchor=west] at (0.8,-1.15) {等快度 $y=\frac12\ln(v^+/v^-)$（对角线）};
    \node[fa cap] at (2.2,-1.85) {(b)\ 对数光锥平面};
  \end{scope}
\end{tikzpicture}
\caption{光锥坐标与 Lorentz 推动。(a) $(z,t)$ 平面中的光锥轴 $x^\pm=(t\pm z)/\sqrt2$ 与沿 $+z$ 快速运动的强子世界带；推动沿不变双曲线移动点。
(b) 对数光锥平面中，快度为 $\eta$ 的推动 $v^\pm\to \ee^{\pm\eta}v^\pm$ 是沿反对角线的平移（$\ln v^+$ 增加 $\eta$，$\ln v^-$ 减少 $\eta$），等 $v^+v^-$ 线为反对角线（虚线），等快度线为对角线（点线）。}
\label{fig:prelim-lightcone}
\end{figure}
%%% END fig:prelim-lightcone
% -------------------------------------------------------------------------------------
\section{Wilson 线}\label{sec:prelim-wilson}

部分子分布（第~8~章）含有位于不同点的场的乘积 $\bar\psi(w)\cdots\psi(0)$；在局域规范变换 $\psi(x)\to U(x)\psi(x)$ 下它变为 $\bar\psi(w)U^\dagger(w)\cdots U(0)\psi(0)$，不是规范不变的，需要一个沿路径的平行移动算符，即 Wilson 线。另一方面，高能夸克吸收软胶子时几乎不改变方向（\cref{fig:soft-eikonal}），在软胶子看来它就是一条沿其速度方向的经典色荷世界线，其与胶子场的耦合也是 Wilson 线。这两重身份在因子化证明中反复出现。

\begin{definition}[Wilson 线]\label{def:prelim-wilson}
(a) 对分段光滑路径 $z(s)$，$s\in[0,1]$，$z(0)=a$，$z(1)=b$，记 $A_\mu\equiv A^a_\mu t^a$，定义
\begin{equation}\label{eq:prelim-W}
\begin{aligned}
W[b,a]&=P\exp\Big(\ii g\int_a^b\dd z^\mu A^a_\mu(z)\,t^a\Big)\\
&\equiv\sum_{m=0}^\infty(\ii g)^m\!\!\int\limits_{0<s_1<\dots<s_m<1}\!\!\dd s_1\cdots\dd s_m\;\dot z\cdot A(z(s_m))\cdots\dot z\cdot A(z(s_1)),
\end{aligned}
\end{equation}
即路径上较晚（离 $b$ 较近）的点排在左边。% VF-CONV-9
\par
(b) 对指向未来的矢量 $n$（$n^0>0$，类光或类时），半无限直线定义为极限 $W[\infty n,y]=\lim_{L\to\infty}W[y+Ln,y]$、$W[y,-\infty n]=\lim_{L\to\infty}W[y,y-Ln]$，以及 $W[y,\infty n]=W[\infty n,y]^\dagger$、$W[-\infty n,y]=W[y,-\infty n]^\dagger$。极限中采用绝热正规化：半无限线上的每个场乘以 $\ee^{-\epsilon s}$，$s\ge0$ 为 $z=y\pm sn$ 中的仿射参数（到有限端点的距离），最后取 $\epsilon\to0^+$。（绝热参数 $\epsilon$ 与维数正规化的 $\varepsilon$ 无关。）
\par
(c) 伴随表示的 Wilson 线 $W^{(A)}_{ab}[b,a]=2\Tr\big[t^aW[b,a]\,t^bW[a,b]\big]$。
\par
(d) 胶子场的 Fourier 约定 $A_\mu(x)=\int\frac{\dd^4k}{(2\pi)^4}\ee^{-\ii k\cdot x}\tilde A_\mu(k)$：$k$ 是在连接点流\emph{入} Wilson 线的动量；厄米性给出 $\tilde A(k)^\dagger=\tilde A(-k)$。
\uses{def:prelim-qcd,def:prelim-lightfront}
\end{definition}

\begin{proposition}[Wilson 线的规范协变性 \tagR]\label{prop:prelim-wilson-covariance}
（$R_{\mathrm{def}}$：由算符定义直接推出。）设胶子场光滑。
\par
(a) $W(s)\equiv W[z(s),a]$ 是初值问题
\begin{equation}\label{eq:prelim-W-ode}
\frac{\dd}{\dd s}W(s)=\ii g\,\dot z\cdot A(z(s))\,W(s),\qquad W(0)=1
\end{equation}
的唯一解；$W$ 是幺正的，行列式为 $1$。
\par
(b) 在\cref{def:prelim-qcd}的规范变换 $A_\mu\to UA_\mu U^{-1}-\frac{\ii}{g}(\partial_\mu U)U^{-1}$ 下，$W[b,a]\to U(b)\,W[b,a]\,U^\dagger(a)$。若把较早的点排在左边（反路径序），则在非阿贝尔情形这一变换律不成立。
\par
(c) 组合与折返：对相接的路径 $W[c,b]W[b,a]=W[c,a]$；沿同一路径折返 $W[a,b]W[b,a]=1$。
\uses{def:prelim-qcd,def:prelim-wilson,prop:qcd-covariance}
\end{proposition}

\begin{proof}
(a) 级数\cref{eq:prelim-W}的第 $m$ 项以 $(g\max_s\|\dot z\cdot A\|)^m/m!$ 为界（单纯形的体积为 $1/m!$），故绝对收敛，逐项求导即得\cref{eq:prelim-W-ode}；线性方程的解唯一（Picard--Lindelöf）。$A$ 厄米给出 $\frac{\dd}{\dd s}(W^\dagger W)=0$，$\Tr t^a=0$ 给出 $\frac{\dd}{\dd s}\det W=\Tr(\ii g\,\dot z\cdot A)\det W=0$。

(b) 令 $V(s)=U(z(s))\,W(s)\,U^\dagger(a)$，则 $V(0)=1$，且
\[
\frac{\dd V}{\dd s}=\Big[\frac{\dd U}{\dd s}U^{-1}+\ii g\,U\,\dot z\cdot A\,U^{-1}\Big]V
=\ii g\,\dot z\cdot\Big[UAU^{-1}-\frac{\ii}{g}(\partial U)U^{-1}\Big]V=\ii g\,\dot z\cdot A'\,V ,
\]
其中用了 $\dot z^\mu\partial_\mu U=\dd U/\dd s$；$V$ 解变换后的场 $A'$ 的初值问题，由唯一性 $V=W'$。

(c) $W[c,b]W[b,a]$ 与 $W[c,a]$ 在第二段路径上满足同一方程、在 $b$ 处初值相同，故相等。对折返路径 $\tilde z(t)=z(1-t)$，$R(t)\equiv W(1-t)W(1)^{-1}$ 满足 $\dd R/\dd t=\ii g\,\dot{\tilde z}\cdot A(\tilde z(t))R$，$R(0)=1$，故 $R(t)=W[\tilde z(t),b]$，取 $t=1$ 得 $W[a,b]=W[b,a]^{-1}$。
\end{proof}

(b) 中反路径序的断言可由无穷小变换 $U=1+\ii\theta$ 验证：$\delta A_\mu=\ii[\theta,A_\mu]+\frac1g\partial_\mu\theta$，对\cref{eq:prelim-W-ode}作变分并分部积分，路径序给出 $\delta W=\ii(\theta(b)W-W\theta(a))$，两个对易子项相消；反路径序（$\dd V/\dd s=\ii g\,V\dot z\cdot A$）则给出 $\delta V=\ii(V\theta(b)-\theta(a)V)+2g\int_0^1\dd s\,V(0,s)[\dot z\cdot A,\theta]V(s,1)$，末项一般不为零。反路径序指数是耦合符号相反的约定 $D=\partial+\ii gA$ 下的协变 Wilson 线（Collins~\cite{Collins2011}的约定，换算时令 $g\to-g$）。在随机的 SU(2)、SU(3) 背景场中对离散化路径序指数的数值检验确认了 (b)、(c) 以及伴随线等于以 $(T^c)_{ab}=-\ii f^{cab}$ 为生成元的路径序指数。

\begin{proposition}[半无限直 Wilson 线的动量空间 Feynman 规则 \tagR]\label{prop:prelim-wilson-rules}
（$R_{\mathrm{def}}$：由算符定义直接推出。）设直路径由起点指向终点的切矢量为 $v$，一端在无穷远，场按\cref{def:prelim-wilson}(d) 展开并带绝热因子。含 $m$ 个胶子（Lorentz 指标 $\mu_j$、色指标 $a_j$、流入动量 $k_j$）的项由下列"统一规则"给出：
\begin{enumerate}[label=(\roman*),leftmargin=2.2em,itemsep=0pt,topsep=2pt]
\item 每个连接点一个顶点因子 $\ii g\,t^{a}v^{\mu}$；
\item 除伸向无穷远的一段外，每一段一个程函传播子 $\ii/(v\cdot\ell+\ii0)$，$\ell$ 是沿路径取向流动的动量，由各连接点的动量守恒与"伸向无穷远的一段 $\ell=0$"确定（连到有限端点的一段携带流入的总动量）；
\item 色矩阵按路径次序排列，靠近终点者在左；
\item 在有限绝热参数下，某一段的 $\ii0$ 等于 $\ii\epsilon$ 乘以该段与无穷远端之间的连接点数。
\end{enumerate}
对指向未来的 $n$ 与 $\int[\dd k]_m\equiv\int\prod_{j=1}^m\dd^4k_j/(2\pi)^4$，$\tilde A\equiv\tilde A^at^a$，胶子按路径次序编号：% VF-WL-1
\begin{align}
W^{(m)}[0,-\infty n]&=\int[\dd k]_m\,\big[\ii g\,n\cdot\tilde A(k_m)\big]\cdots\big[\ii g\,n\cdot\tilde A(k_1)\big]\prod_{j=1}^m\frac{\ii}{n\cdot(k_1+\dots+k_j)+\ii0},\label{eq:prelim-past}\\
W^{(m)}[\infty n,0]&=\int[\dd k]_m\,\big[\ii g\,n\cdot\tilde A(k_m)\big]\cdots\big[\ii g\,n\cdot\tilde A(k_1)\big]\prod_{j=1}^m\frac{-\ii}{n\cdot(k_j+\dots+k_m)-\ii0},\label{eq:prelim-future}
\end{align}
在\cref{eq:prelim-past}中 $k_1$ 最靠近 $-\infty$，在\cref{eq:prelim-future}中 $k_1$ 最靠近 $0$。取向相反的 $W[-\infty n,0]$ 与 $W[0,\infty n]$ 的规则由上两式取厄米共轭得到（顶点变为 $-\ii g\,t^an^\mu$，传播子取复共轭，并把 $k_j$ 换成 $-k_j$），也可以直接用 (i)--(iv) 读出（此时 $v=-n$）；割图中割线右侧（共轭振幅）的 Wilson 线同样使用复共轭的规则。特别地，单胶子项为
\begin{equation}\label{eq:prelim-one-gluon}
W^{(1)}[\infty n,0]=g\,t^a\,\frac{n\cdot\tilde A^a(k)}{n\cdot k-\ii0},\qquad W^{(1)}[-\infty n,0]=g\,t^a\,\frac{n\cdot\tilde A^a(k)}{n\cdot k+\ii0},
\end{equation}
二者之差（同一端点处"指向未来减指向过去"）为 $2\pi\ii\,g\,t^a\,n\cdot\tilde A^a(k)\,\delta(n\cdot k)$。
\uses{def:prelim-wilson,def:prelim-feynman-rules,prop:appB-rules}
\end{proposition}

\begin{proof}
完整证明见附录~B（\cref{prop:appB-rules}）；要点如下。$W[0,-\infty n]$ 的第 $m$ 项中代入 Fourier 展开与绝热因子 $\ee^{\epsilon\lambda_j}$，参数积分为 $\int_{-\infty<\lambda_1<\dots<\lambda_m<0}\prod_j\dd\lambda_j\,\ee^{-\ii\lambda_j(n\cdot k_j+\ii\epsilon)}$。先积 $\lambda_1$ 得 $\frac{\ii}{n\cdot k_1+\ii\epsilon}\ee^{-\ii\lambda_2(n\cdot k_1+\ii\epsilon)}$，它与第二个场的相因子合并；归纳地得到 $\prod_j\ii/[n\cdot(k_1+\dots+k_j)+\ii j\epsilon]$，色矩阵 $t^{a_m}\cdots t^{a_1}$ 的次序即路径次序：这就是\cref{eq:prelim-past}与 (iv)。对 $W[\infty n,0]$，区域为 $0<\lambda_1<\dots<\lambda_m<\infty$，阻尼 $\ee^{-\epsilon\lambda_j}$，先积 $\lambda_m$ 即得\cref{eq:prelim-future}。取向相反的线由厄米共轭与 $\tilde A(k)^\dagger=\tilde A(-k)$ 得到；最后 $\frac1{x-\ii0}-\frac1{x+\ii0}=2\pi\ii\,\delta(x)$。
\end{proof}

\cref{fig:prelim-wilson-rules}把\cref{eq:prelim-past,eq:prelim-future}画成图。规则在 $n\to\alpha n$（$\alpha>0$）下不变，但 $\alpha<0$ 时 $\ii0$ 反号，所以 $n$ 必须沿部分子的物理运动方向；"$+\ii0$ 配沿取向流动的动量"正是被替代的 Feynman 传播子的程函极限（\cref{ex:prelim-eikonal-two}）。

%%% BEGIN fig:prelim-wilson-rules
\begin{figure}[htbp]
\centering
\begin{tikzpicture}[fd,
  fa end/.style={circle,draw,fill=white,line width=0.8pt,inner sep=0pt,minimum size=5.5pt},
  fa idx/.style={font=\footnotesize,inner sep=1.5pt},
  fa prop/.style={font=\small,inner sep=2pt},
  fa cap/.style={font=\footnotesize}]
  % ===================== (a) past-pointing line W[0,-infinity n] =====================
  \begin{scope}
    \draw[wilson=0.06] (0,0) -- (12.3,0);
    \node[fa idx,left] at (-0.05,0) {$-\infty n$};
    \node[fa end] at (12.3,0) {};
    \node[fa idx,right] at (12.4,0) {$0$};
    \foreach \x/\j in {1.5/1,5.1/2,8.7/3} {
      \draw[gluon] (\x,1.35) node[fa idx,above]{$\mu_\j,a_\j$} -- (\x,0);
      \path[mom={$k_\j$},mom sep=0.3cm] (\x,1.35) -- (\x,0.12);
      \node[vertex] at (\x,0) {};
    }
    \node[fa prop,below] at (0.75,-0.18) {$\ell=0$};
    \node[fa idx,anchor=south east] at (1.38,0.13) {$\ii g\,t^{a_1}n^{\mu_1}$};
    \node[fa prop,below] at (3.3,-0.1) {$\dfrac{\ii}{n\cdot k_1+\ii0}$};
    \node[fa prop,below] at (6.9,-0.1) {$\dfrac{\ii}{n\cdot(k_1+k_2)+\ii0}$};
    \node[fa prop,below] at (10.5,-0.1) {$\dfrac{\ii}{n\cdot(k_1+k_2+k_3)+\ii0}$};
    \node[fa cap] at (6.15,-1.55) {(a)\ 过去指向的 $W[0,-\infty n]$：$k_1$ 最靠近 $-\infty$，颜色因子 $t^{a_3}t^{a_2}t^{a_1}$};
  \end{scope}
  % ===================== (b) future-pointing line W[infinity n,0] =====================
  \begin{scope}[yshift=-4.1cm]
    \draw[wilson=0.94] (0,0) -- (12.3,0);
    \node[fa end] at (0,0) {};
    \node[fa idx,left] at (-0.1,0) {$0$};
    \node[fa idx,right] at (12.35,0) {$+\infty n$};
    \foreach \x/\j in {3.6/1,7.2/2,10.8/3} {
      \draw[gluon] (\x,1.35) node[fa idx,above]{$\mu_\j,a_\j$} -- (\x,0);
      \path[mom={$k_\j$},mom sep=0.3cm] (\x,1.35) -- (\x,0.12);
      \node[vertex] at (\x,0) {};
    }
    \node[fa idx,anchor=south east] at (3.48,0.13) {$\ii g\,t^{a_1}n^{\mu_1}$};
    \node[fa prop,below] at (1.8,-0.1) {$\dfrac{-\ii}{n\cdot(k_1+k_2+k_3)-\ii0}$};
    \node[fa prop,below] at (5.4,-0.1) {$\dfrac{-\ii}{n\cdot(k_2+k_3)-\ii0}$};
    \node[fa prop,below] at (9.0,-0.1) {$\dfrac{-\ii}{n\cdot k_3-\ii0}$};
    \node[fa prop,below] at (11.55,-0.18) {$\ell=0$};
    \node[fa cap] at (6.15,-1.55) {(b)\ 未来指向的 $W[\infty n,0]$：$k_1$ 最靠近 $0$，颜色因子 $t^{a_3}t^{a_2}t^{a_1}$};
  \end{scope}
\end{tikzpicture}
\caption{半无限 Wilson 线的动量空间规则（$k$ 流入 Wilson 线；颜色矩阵按路径排序，靠近终点者在左）：(a) 过去指向，(b) 未来指向；双线上的箭头为路径取向，空心圆为有限端点 $0$。
统一规则（VF-WL-1）：每个顶点 $\ii g\,t^{a_j}n^{\mu_j}$（图中只标出第一个）；有限端点与顶点之间、相邻顶点之间的每一段给一个程函传播子 $\ii/(n\cdot\ell+\ii0)$，
$\ell$ 为沿路径取向流动的动量，伸向无穷远的一段 $\ell=0$、不给传播子。于是 (a) 中第 $j$ 个顶点之后的一段 $\ell=k_1+\dots+k_j$；
(b) 中第 $j$ 个顶点之前的一段 $\ell=-(k_j+\dots+k_3)$，传播子即 $-\ii/(n\cdot(k_j+\dots+k_3)-\ii0)$。}
\label{fig:prelim-wilson-rules}
\end{figure}
%%% END fig:prelim-wilson-rules

对连接次序求和时，若颜色矩阵彼此对易（阿贝尔情形或颜色对称化的部分），嵌套的程函分母合成单胶子因子之积：

\begin{lemma}[程函恒等式]\label{lem:prelim-eikonal-identity}
设复数 $a_1,\dots,a_m$ 的一切部分和 $a_{i_1}+\dots+a_{i_r}$（$\{i_1,\dots,i_r\}$ 取遍 $\{1,\dots,m\}$ 的非空子集）都不为零，则
\begin{equation}\label{eq:prelim-eikonal}
\sum_{\pi\in S_m}\prod_{j=1}^m\frac1{a_{\pi(1)}+\dots+a_{\pi(j)}}=\prod_{j=1}^m\frac1{a_j} .
\end{equation}
特别地，取 $a_j=n\cdot k_j+\ii\epsilon$（$\epsilon>0$），左边第 $j$ 个分母为 $n\cdot(k_{\pi(1)}+\dots+k_{\pi(j)})+\ii j\epsilon$：带绝热正规化的阿贝尔程函恒等式在有限 $\epsilon$ 下严格成立。若在每个分母上都加同一个 $+\ii\epsilon$，则恒等式在有限 $\epsilon$ 下不成立，只在 $\epsilon\to0$ 时作为分布成立。% VF-WL-2
\uses{lem:appB-shuffle}
\end{lemma}

\begin{proof}
（单纯形分解）若一切 $\Imag a_j>0$，则 $\prod_j\ii/a_j=\prod_j\int_{-\infty}^0\dd\lambda_j\,\ee^{-\ii\lambda_ja_j}$；$(-\infty,0]^m$ 除零测集外是 $m!$ 个单纯形 $\{\lambda_{\pi(1)}<\dots<\lambda_{\pi(m)}\}$ 的不交并，与\cref{eq:prelim-past}的推导相同的逐次积分表明，每个单纯形贡献 $\prod_j\ii/(a_{\pi(1)}+\dots+a_{\pi(j)})$。两边是有理函数，在开集上相等，故处处相等（物理上：阿贝尔 Wilson 线是普通指数，$W^{(m)}=(W^{(1)})^m/m!$）。

（归纳）按 $\pi(m)=l$ 把置换分组：末因子 $1/(a_1+\dots+a_m)$ 公共，其余因子对 $\{1,\dots,m\}\setminus\{l\}$ 的置换求和，由归纳假设等于 $\prod_{j\ne l}1/a_j$；于是左边 $=\frac1{a_1+\dots+a_m}\sum_l\prod_{j\ne l}\frac1{a_j}=\prod_j\frac1{a_j}$。

取 $a_j=n\cdot k_j+\ii\epsilon$ 时部分和的虚部为 $r\epsilon>0$。若每个分母都加 $+\ii\epsilon$（$a_j=n\cdot k_j$ 实），则 $m=2$ 时
\[
\sum_{\pi\in S_2}\prod_{j=1}^2\frac{\ii}{a_{\pi(1)}+\dots+a_{\pi(j)}+\ii\epsilon}-\prod_{j=1}^2\frac{\ii}{a_j+\ii\epsilon}=\frac{-\ii\epsilon}{(a_1+\ii\epsilon)(a_2+\ii\epsilon)(a_1+a_2+\ii\epsilon)}\ne0 ;
\]
它与检验函数 $[(1+a_1^2)(1+a_2^2)]^{-1}$ 的积分为 $\pi^2\epsilon/[(1+\epsilon)^2(2+\epsilon)]$，只在 $\epsilon\to0$ 时为零。洗牌形式见\cref{lem:appB-shuffle}。
\end{proof}

阿贝尔程函恒等式本身是经典的。统一的 $+\ii\epsilon$ 相当于只对离无穷远端最远的连接点加阻尼，并不是 Wilson 线算符的正规化；算符层次的指数阻尼给出累积分母（第 $j$ 个分母带 $j$ 倍正规化参数）并保持指数化，统一替换则破坏指数化，见 Echevarria、Scimemi 与 Vladimirov~\cite[式~(5)--(6) 及其后的讨论]{Echevarria:2015byo}。非阿贝尔情形中颜色按路径排序，\cref{eq:prelim-eikonal}只适用于颜色对称化的部分；反对称部分由对易子 $[t^b,t^a]=\ii f^{bac}t^c$ 组成，即携带总动量的伴随"网"。例如两个胶子时 $W^{(2)}-\frac12(W^{(1)})^2$ 只含 $[T_2,T_1]$ 乘以 $E(1,2)-E(2,1)$（$E(r,s)$ 为先连接 $r$ 后连接 $s$ 的程函传播子之积），这由\cref{eq:prelim-eikonal}把 $E(1)E(2)$ 写成 $E(1,2)+E(2,1)$ 即得。% VF-WL-3
以网表述的全阶非阿贝尔指数化是全阶论证（\cref{prop:appB-webs}；Gatheral~\cite{Gatheral1983}，Frenkel 与 Taylor~\cite{FrenkelTaylor1984}）。

\begin{example}[$m=2$ 的程函恒等式与入射夸克传播子的程函极限]\label{ex:prelim-eikonal-two}
(a) $m=2$ 时\cref{eq:prelim-eikonal}即 $\frac1{a_1(a_1+a_2)}+\frac1{a_2(a_1+a_2)}=\frac{a_1+a_2}{a_1a_2(a_1+a_2)}=\frac1{a_1a_2}$；带绝热正规化时，$\frac{\ii}{a_1+\ii\epsilon}\frac{\ii}{a_1+a_2+2\ii\epsilon}+\frac{\ii}{a_2+\ii\epsilon}\frac{\ii}{a_1+a_2+2\ii\epsilon}$ 的分子之和 $(a_1+\ii\epsilon)+(a_2+\ii\epsilon)$ 恰好约去 $a_1+a_2+2\ii\epsilon$，得 $\frac{\ii}{a_1+\ii\epsilon}\frac{\ii}{a_2+\ii\epsilon}$。
\par
(b) 一个动量为 $p$（$p^2=0$）的入射夸克吸收一个软胶子（动量 $k$ 流入夸克线，色 $a$，指标 $\mu$；$k$ 的各分量都远小于 $p$ 的大分量，且不在 Glauber 区）。由 $\sla pu(p)=0$，$(\sla p+\sla k)\gamma^\mu u(p)=2p^\mu u(p)+\sla k\gamma^\mu u(p)$，$(p+k)^2=2p\cdot k+k^2$。略去相对为 $O(k)$ 的 $\sla k\gamma^\mu$ 与相对于 $p\cdot k$ 的 $k^2$，
\begin{equation}\label{eq:prelim-eikonal-in}
S(p+k)\,\ii g\gamma^\mu t^a\,u(p)=\frac{\ii(\sla p+\sla k)}{(p+k)^2+\ii0}\,\ii g\gamma^\mu t^a\,u(p)\simeq-g\,t^a\,\frac{p^\mu}{p\cdot k+\ii0}\,u(p). % VF-WL-5
\end{equation}
它与过去指向的 Wilson 线 $W[0,-\infty n]$（$n\propto p$）的单胶子规则 $(\ii g\,t^an^\mu)\,\ii/(n\cdot k+\ii0)=-g\,t^an^\mu/(n\cdot k+\ii0)$ 完全相同（规则在 $n\to\alpha n$，$\alpha>0$ 下不变）。同样，出射夸克给出 $\bar u(p)\,\ii g\gamma^\mu t^a S(p-k)\simeq+g\,t^a\,p^\mu\bar u(p)/(p\cdot k-\ii0)$，即未来指向的 $W[\infty n,0]$ 的规则 $(\ii g\,t^an^\mu)(-\ii)/(n\cdot k-\ii0)$。$\ii0$ 的符号就是被替代的传播子 $(p\pm k)^2+\ii0$ 的符号。
\par
(c) 两个软胶子依次（先 $k_1$ 后 $k_2$）被吸收时，重复 (b) 两次得\cref{eq:prelim-past}中这一次序的项 $\frac{-g\,p^{\mu_2}}{p\cdot(k_1+k_2)+\ii0}\frac{-g\,p^{\mu_1}}{p\cdot k_1+\ii0}t^{a_2}t^{a_1}u(p)$；对两种次序求和，颜色对易时由 (a) 得到单胶子因子之积，QCD 中则得到路径序的 $W^{(2)}$。
\uses{lem:prelim-eikonal-identity,prop:prelim-wilson-rules}
\end{example}

\cref{fig:soft-eikonal}画出\cref{ex:prelim-eikonal-two}(b) 的两种情形；图中的 $n$ 是夸克动量 $p$ 的方向（不是本书的 $n_A,n_B$），图注所说的 Glauber 区见\cref{def:ir-glauber}（第~4~章）。程函近似的相对误差是 $O(\lambda)$（$k$ 整体缩小为 $\lambda$ 倍），但系数正比于 $1/|p\cdot k|$：在 $p\cdot k\to0$ 而 $k^2\ne0$ 的 Glauber 区近似失效。

%%% BEGIN fig:soft-eikonal
\begin{figure}[htbp]
\centering
\begin{tikzpicture}[fd,
  t2 soft/.style={draw=green!45!black},
  t2 lab/.style={font=\small},
  t2 eq/.style={font=\large},
  t2 hc/.style={hard blob,circle,minimum size=0.84cm,inner sep=0pt}]
  % ======================= (a) outgoing quark =======================
  \node[t2 hc] (H) at (0,0) {$H$};
  \coordinate (X) at (1.95,0);
  \draw[quark] (H.east) -- (X);
  \draw[quark] (X) -- (3.35,0);
  \draw[gluon,t2 soft] (2.6,1.3) -- (X);
  \node[vertex] at (X) {};
  \path[mom'={$p-k$}] (H.east) -- (X);
  \path[mom'={$p$}] (X) -- (3.35,0);
  \path[mom'={$k$},mom sep=0.3cm] (2.6,1.3) -- (X);
  \node[t2 eq] at (4.15,0) {$\approx$};
  \begin{scope}[xshift=5.4cm]
    \node[t2 hc] (H) at (0,0) {$H$};
    \draw[quark] (H.340) -- ($(H.340)+(3.0,0)$) coordinate (Eq);
    \draw[wilson=0.82] (H.20) -- ($(H.20)+(3.0,0)$) coordinate (Ew);
    \coordinate (Y) at ($(H.20)!0.55!(Ew)$);
    \draw[gluon,t2 soft] ($(Y)+(0.65,1.16)$) -- (Y);
    \path[mom'={$p$},mom pos=0.75] (H.340) -- (Eq);
    \node[t2 lab,anchor=south] at ($(H.20)!0.25!(Ew)+(0,0.1)$) {$W_n[\infty,0]$};
  \end{scope}
  \node[font=\small] at (1.45,-1.15) {$\bar u(p)\,(\ii g\,t^a\gamma^\mu)\,\dfrac{\ii(\sla p-\sla k)}{(p-k)^2+\ii0}$};
  \node[font=\small] at (7.0,-1.15) {$\bar u(p)\,\dfrac{g\,t^a\,n^\mu}{n\cdot k-\ii0}$};
  \node[font=\small,anchor=west] at (-0.6,-2.0) {(a)\ 出射夸克：未来指向的 Wilson 线（$n\propto p$）};
  % ======================= (b) incoming quark =======================
  \begin{scope}[yshift=-4.1cm]
    \node[t2 hc] (H) at (3.35,0) {$H$};
    \coordinate (X) at (1.4,0);
    \draw[quark] (0,0) -- (X);
    \draw[quark] (X) -- (H.west);
    \draw[gluon,t2 soft] (0.75,1.3) -- (X);
    \node[vertex] at (X) {};
    \path[mom'={$p$}] (0,0) -- (X);
    \path[mom'={$p+k$}] (X) -- (H.west);
    \path[mom={$k$},mom sep=0.3cm] (0.75,1.3) -- (X);
    \node[t2 eq] at (4.15,0) {$\approx$};
    \begin{scope}[xshift=5.4cm]
      \node[t2 hc] (H) at (3.35,0) {$H$};
      \draw[quark] ($(H.200)-(3.0,0)$) coordinate (Sq) -- (H.200);
      \draw[wilson=0.25] ($(H.160)-(3.0,0)$) coordinate (Sw) -- (H.160);
      \coordinate (Y) at ($(Sw)!0.45!(H.160)$);
      \draw[gluon,t2 soft] ($(Y)+(-0.65,1.16)$) -- (Y);
      \path[mom'={$p$},mom pos=0.3] (Sq) -- (H.200);
      \node[t2 lab,anchor=south] at ($(Y)!0.55!(H.160)+(0,0.1)$) {$W_n[0,-\infty]$};
    \end{scope}
    \node[font=\small] at (1.7,-1.15) {$\dfrac{\ii(\sla p+\sla k)}{(p+k)^2+\ii0}\,(\ii g\,t^a\gamma^\mu)\,u(p)$};
    \node[font=\small] at (7.0,-1.15) {$-\dfrac{g\,t^a\,n^\mu}{n\cdot k+\ii0}\,u(p)$};
    \node[font=\small,anchor=west] at (-0.6,-2.0) {(b)\ 入射夸克：过去指向的 Wilson 线（$n\propto p$）};
  \end{scope}
\end{tikzpicture}
\caption{软胶子的程函（eikonal）近似。软胶子（绿色，动量 $k$ 流入夸克线，色指标 $a$、Lorentz 指标 $\mu$）附着在高能无质量夸克线（动量 $p$）上。
当 $k$ 的各分量都远小于 $p$ 的大分量、且不在 Glauber 区时，分子 $\bar u(p)\gamma^\mu(\sla p-\sla k)\to2p^\mu\bar u(p)$、分母 $(p-k)^2+\ii0\to-2p\cdot k+\ii0$，
结果只依赖于 $p$ 的方向 $n$（对 $n$ 的标度不变），正是沿 $n$ 方向的 Wilson 线的单胶子 Feynman 规则：
(a)~出射夸克对应未来指向的 $W_n[\infty,0]$，传播子为 $-\ii/(n\cdot k-\ii0)$；(b)~入射夸克对应过去指向的 $W_n[0,-\infty]$，传播子为 $\ii/(n\cdot k+\ii0)$
（顶点因子均为 $\ii g\,t^a n^\mu$）。Wilson 线上的箭头与被替代的夸克线的费米子数流向一致。}
\label{fig:soft-eikonal}
\end{figure}
%%% END fig:soft-eikonal

\begin{example}[有限类光线段的单胶子项]\label{ex:prelim-wilson-segment}
取 $w=(0^+,w^-,\mathbf 0_{\mathrm T})=w^-n_B$ 与沿 $n_B$ 的直线段 $z(s)=s\,n_B$，$0\le s\le w^-$。由于 $\dd z\cdot A=\dd s\,n_B\cdot A=\dd s\,A^+$，
\[
W[w,0]=P\exp\Big(\ii g\int_0^{w^-}\dd s\,A^{a+}(s\,n_B)\,t^a\Big).
\]
代入 Fourier 展开并用 $k\cdot n_B=k^+$，单胶子项为
\begin{equation}\label{eq:prelim-segment}
W^{(1)}[w,0]=\ii g\,t^a\!\int\!\frac{\dd^4k}{(2\pi)^4}\,\tilde A^{a+}(k)\int_0^{w^-}\!\dd s\,\ee^{-\ii sk^+}=g\,t^a\!\int\!\frac{\dd^4k}{(2\pi)^4}\,\tilde A^{a+}(k)\,\frac{1-\ee^{-\ii k^+w^-}}{k^+}.
\end{equation}
(i) $k^+\to0$ 时 $(1-\ee^{-\ii k^+w^-})/k^+\to\ii w^-$：极点是可去的，不需要 $\ii0$。
(ii)(iii) 令 $w^-\to\pm\infty$ 并带阻尼 $\ee^{\mp\epsilon s}$：$\int_0^{\pm\infty}\dd s\,\ee^{-\ii sk^+\mp\epsilon s}=-\ii/(k^+\mp\ii\epsilon)$，于是 $W^{(1)}\to g\,t^a\tilde A^{a+}(k)/(k^+\mp\ii0)$，分别是\cref{eq:prelim-one-gluon}中 $W^{(1)}[\pm\infty n_B,0]$ 的规则。
(iv) 经由 $+\infty$ 分解：$W^{(1)}[w,\infty n_B]=-g\,t^a\ee^{-\ii k^+w^-}\tilde A^{a+}/(k^+-\ii0)$（共轭规则加端点相因子），与 (ii) 相加得 $g\,t^a\tilde A^{a+}(1-\ee^{-\ii k^+w^-})/(k^+-\ii0)$，因分子在 $k^+=0$ 处为零而等于\cref{eq:prelim-segment}；经由 $-\infty$ 同理。所以对同一条类光线上的两点，伸向未来还是过去无关紧要：
\[
W[w,0]=W[w,\infty n_B]\,W[\infty n_B,0]=W[w,-\infty n_B]\,W[-\infty n_B,0].
\]
\uses{def:prelim-wilson,prop:prelim-wilson-rules}
\end{example}

\cref{eq:prelim-segment}是夸克部分子分布（\cref{def:pdf-quark}，第~8~章）中 Wilson 线的单胶子项，附录~F 的伸缩求和例子也从它出发；对横动量依赖的分布，两端点不在同一条类光线上，取向就不再无关（第~14~章）。

类光 Wilson 线还有两种在后文至关重要的发散。其一，两条直 Wilson 线在一点相交（"尖点"）时有紫外的尖点发散，其反常维度在类光极限下随两线之间的快度线性增长~\cite{Polyakov1980,BrandtNeriSato1981,KorchemskyRadyushkin1987}，这是 Sudakov 双对数的来源（\cref{def:sudakov-cusp}，第~5~章）。其二，类光线在 Feynman 规范中没有自能（正比于 $n^2=0$），但程函传播子使 $\int\dd k^+/k^+$ 在 $k^+\to0$ 处发散，维数正规化控制不了这一端点：在共线分布中它在固定纵向动量分数时于实、虚修正之间相消，形成加号分布（第~8~章）；在固定横动量的分布中不相消，成为快度发散，需要快度正规子与软因子（第~14~章）。
% -------------------------------------------------------------------------------------
\section{维数正规化的技术约定与 $\gamma_5$}\label{sec:prelim-dimreg}

维数正规化（\cref{def:qcd-dimreg}）保持规范与 BRST 对称性，并同时正规化紫外与红外发散。后一点带来一个陷阱：同一个 $1/\varepsilon$ 可以来自圈动量趋于无穷的区域，也可以来自软、共线区域，"无标度"的积分则因二者相消而为零。因子化把红外极点吸收进部分子分布、把紫外极点交给重整化，所以每个极点的来源都必须记清楚。

\begin{definition}[维数正规化的技术约定]\label{def:prelim-dr-conventions}
本书的维数正规化以\cref{def:qcd-dimreg}（$D=4-2\varepsilon$）与\cref{def:prelim-msbar}（$\MSbar$，$S_\varepsilon=(4\pi)^\varepsilon\ee^{-\varepsilon\gamma_E}$）为基础，另作以下约定。
\par
(a) 极点来源的标注：$1/\varepsilon_{\mathrm{UV}}$ 来自圈动量趋于无穷的区域，由重整化抵消项消去；$1/\varepsilon_{\mathrm{IR}}$ 来自软或共线区域（零动量或无质量线彼此共线的位形），在红外安全量中相消或被吸收进强子矩阵元。二者是同一个参数 $\varepsilon$，下标只是记账：紫外部分在 $\varepsilon>0$ 处收敛，红外部分在 $\varepsilon<0$ 处收敛，各自从收敛区解析延拓。
\par
(b) 初态平均：夸克自旋取 $\frac12$、颜色取 $\frac1{N_c}$；胶子极化取 $\frac1{D-2}=\frac1{2(1-\varepsilon)}$、颜色取 $\frac1{N_c^2-1}$。胶子极化的 $D$ 维平均是部分子系数函数 $\MSbar$ 约定的一部分：若改为对 $2$ 个极化平均，胶子道的系数函数平移 $C_2^g\to C_2^g+P_{qg}$，$\Delta_{qg}\to\Delta_{qg}+P_{qg}$（第~10~章与第~12~章的记号）。% VF-CONV-6
\par
(c) Dirac 代数用常规维数正规化（CDR）：$\{\gamma^\mu,\gamma^\nu\}=2g^{\mu\nu}$，$g^\mu{}_\mu=D$，$\Tr\mathbb 1=4$，所有动量与极化矢量都在 $D$ 维中；$D$ 维相空间以及 $S_\varepsilon$ 约定的影响见\cref{sec:appC-msbar}，$\gamma_5$ 的约定见本节末。
\uses{def:qcd-dimreg,def:prelim-msbar}
\end{definition}

(b) 的理由：$D$ 维中横向极化有 $D-2$ 个；若对 $2$ 个极化平均，胶子道部分子截面乘以 $1-\varepsilon$，与共线极点 $-\frac2\varepsilon P_{qg}$ 相乘留下有限项 $2P_{qg}$，即上述平移。(c) 中与 Collins~\cite{Collins2011}的 $S_\varepsilon^{\mathrm C}=(4\pi)^\varepsilon/\Gamma(1-\varepsilon)$ 相比，$S_\varepsilon/S^{\mathrm C}_\varepsilon=1+\frac{\pi^2}{12}\varepsilon^2+O(\varepsilon^3)$，只在扣除 $1/\varepsilon^2$ 极点处改变有限常数。% VF-CONV-4

\begin{lemma}[无标度积分为零]\label{lem:prelim-scaleless}
(a) 对任意复数 $a$，在维数正规化中
\begin{equation}\label{eq:prelim-scaleless}
\int\dd^Dk\,(-k^2-\ii0)^{-a}=0 .
\end{equation}
确切地说：作 Wick 转动后，把径向积分在欧氏模 $|k_E|=\Lambda$ 处分开，红外部分在 $\Real D>2\Real a$ 处收敛，紫外部分在 $\Real D<2\Real a$ 处收敛；二者各自解析延拓为 $D$ 的亚纯函数，其和恒为零且与 $\Lambda$ 无关。对 $a=2$、$D=4-2\varepsilon$，
\begin{equation}\label{eq:prelim-scaleless-a2}
\int\dd^Dk\,(-k^2-\ii0)^{-2}=\frac{\ii\pi^{D/2}}{\Gamma(D/2)}\,\Lambda^{-2\varepsilon}\Big(\frac1{\varepsilon_{\mathrm{UV}}}-\frac1{\varepsilon_{\mathrm{IR}}}\Big)=0 .
\end{equation}
\par
(b) 更一般地，被积函数不含任何质量或外部不变量、因而在全部圈动量的整体伸缩下齐次的维数正规化积分为零。
\uses{def:qcd-dimreg}
\end{lemma}

\begin{proof}
(a) 被积函数的奇点在 $k^0=\pm(|\mathbf k|-\ii0)$（第四与第二象限），把 $k^0$ 的路径逆时针转到虚轴（$k^0=\ii k_4$）不越过奇点，故积分等于 $\ii\int\dd^Dk_E\,(k_E^2)^{-a}$。用 $\dd^Dk_E=\Omega_DK^{D-1}\dd K$，$\Omega_D=2\pi^{D/2}/\Gamma(D/2)$：
\begin{gather*}
\int_0^\Lambda\dd K\,K^{D-1-2a}=\frac{\Lambda^{D-2a}}{D-2a}\quad(\Real D>2\Real a),\\
\int_\Lambda^\infty\dd K\,K^{D-1-2a}=-\frac{\Lambda^{D-2a}}{D-2a}\quad(\Real D<2\Real a).
\end{gather*}
右端都是 $D$ 的亚纯函数，给出两部分的解析延拓，其和恒为零且与 $\Lambda$ 无关。这个积分没有共同的收敛区，维数正规化把它定义为两部分延拓之和，与维数正规化的公理（线性、平移不变、伸缩性质；Collins~\cite[第~4.1~节]{Collins1984}）给出的唯一值一致。$a=2$ 时红外部分为 $-\Lambda^{-2\varepsilon}/(2\varepsilon_{\mathrm{IR}})$，紫外部分为 $+\Lambda^{-2\varepsilon}/(2\varepsilon_{\mathrm{UV}})$，乘以 $\ii\Omega_D$ 即得\cref{eq:prelim-scaleless-a2}。

(b) 维数正规化积分满足伸缩性质 $\int\dd^Dk\,f(\rho k)=\rho^{-D}\int\dd^Dk\,f(k)$（对收敛积分即变量替换，一般情形是定义性公理之一）。若 $f(\rho k)=\rho^{-2a}f(k)$，则 $(\rho^{-2a}-\rho^{-D})\int f=0$ 对一切 $\rho>0$ 成立，故 $\int f\equiv0$；多圈积分对全部圈动量同时伸缩，论证相同。
\end{proof}

无标度的形式有时要经过一步变换才显出来：$p^2=0$ 时 $\int\dd^Dl\,[l^2(l+p)^2]^{-1}$ 经 Feynman 参数化 $x\,l^2+(1-x)(l+p)^2=(l+(1-x)p)^2+x(1-x)p^2$ 并平移圈动量后只依赖于 $l'^2$，由 (a) 为零。

\begin{example}[无标度积分的紫外/红外记账]\label{ex:prelim-scaleless}
(i) \emph{在壳无质量夸克的自能。}Feynman 规范中单圈自能 $-\ii\Sigma(p)$ 的计算用到 $\gamma^\mu\sla a\gamma_\mu=-(D-2)\sla a$、对称性 $\int\dd^Dl\,l^\mu/[l^2(p-l)^2]=\frac12p^\mu\int\dd^Dl\,[l^2(p-l)^2]^{-1}$ 与附录~C 的两点积分，结果为
\begin{equation}\label{eq:prelim-sigma}
\begin{aligned}
\Sigma(p)&=\sla p\,\Sigma_V(p^2),\qquad \Sigma_V(p^2)=-\frac{\as C_F}{4\pi}\,(1-\varepsilon)\,\ee^{\varepsilon\gamma_E}\frac{\Gamma(\varepsilon)\Gamma(1-\varepsilon)^2}{\Gamma(2-2\varepsilon)}\Big(\frac{\mu^2}{-p^2-\ii0}\Big)^{\varepsilon}\\
&\phantom{=\sla p\,\Sigma_V(p^2),\qquad \Sigma_V(p^2)}=-\frac{\as C_F}{4\pi}\Big[\frac1\varepsilon+1+\ln\frac{\mu^2}{-p^2-\ii0}+O(\varepsilon)\Big]
\end{aligned}
\end{equation}
（已用 $\MSbar$ 的 $g_0^2=4\pi\as\mu^{2\varepsilon}S_\varepsilon^{-1}$）。它的紫外极点 $-\frac{\as C_F}{4\pi\varepsilon}$ 正是 Feynman 规范的 $\MSbar$ 常数 $\delta Z_2=-\frac{\as C_F}{4\pi\varepsilon}$。在壳（$p^2\to0$）时 $\Sigma_V(0)$ 是无标度积分，由\cref{lem:prelim-scaleless}为零；按来源记账，$\Sigma_V(0)=-\frac{\as C_F}{4\pi}\big(\frac1{\varepsilon_{\mathrm{UV}}}-\frac1{\varepsilon_{\mathrm{IR}}}\big)$。因此裸传播子 $\ii/[\sla p(1-\Sigma_V)]$ 在 $p^2=0$ 处的留数为 $1$：在维数正规化中 $Z_2^{\mathrm{OS}}=1$（每一阶的在壳修正都是无标度的），用裸场写出的振幅外腿不带任何修正。若改用 $\MSbar$ 重整化的场，外腿留数为 $Z_2^{\mathrm{OS}}/Z_2=1+\frac{\as C_F}{4\pi\varepsilon}$：原来的紫外抵消项现在代表一个红外极点。Sudakov 形状因子（\cref{prop:sudakov-oneloop}，第~5~章）的全部极点因此都是红外的。
\par
(ii) \emph{夸克中的夸克分布。}无质量在壳夸克的单圈实辐射图正比于
\begin{equation}\label{eq:prelim-lT-scaleless}
\int\frac{\dd^{2-2\varepsilon}\mathbf l_{\mathrm T}}{\mathbf l_{\mathrm T}^2}
=\underbrace{\int\frac{\dd^{2-2\varepsilon}\mathbf l_{\mathrm T}}{\mathbf l_{\mathrm T}^2+M^2}}_{\pi^{1-\varepsilon}\Gamma(1+\varepsilon)M^{-2\varepsilon}/\varepsilon_{\mathrm{UV}}}
+\underbrace{\int\dd^{2-2\varepsilon}\mathbf l_{\mathrm T}\,\frac{M^2}{\mathbf l_{\mathrm T}^2(\mathbf l_{\mathrm T}^2+M^2)}}_{-\pi^{1-\varepsilon}\Gamma(1+\varepsilon)M^{-2\varepsilon}/\varepsilon_{\mathrm{IR}}}
=\pi^{1-\varepsilon}\Gamma(1+\varepsilon)M^{-2\varepsilon}\Big(\frac1{\varepsilon_{\mathrm{UV}}}-\frac1{\varepsilon_{\mathrm{IR}}}\Big)=0, % VF-PDF-5
\end{equation}
其中辅助标度 $M$ 任意（第一个积分在 $\varepsilon>0$ 处收敛，第二个在 $\varepsilon<0$ 处收敛），虚修正图同样无标度。所以裸分布 $f^{\mathrm{bare}}_{q/q}(\xi)=\delta(1-\xi)$ 对一切阶成立；$\MSbar$ 重整化只减去紫外极点，留下
\[
f_{q/q}(\xi;\mu)=\delta(1-\xi)-\frac{\as}{2\pi}\frac1\varepsilon P_{qq}(\xi)+O(\as^2),
\]
其中剩下的 $1/\varepsilon$ 是红外极点（\cref{prop:pdf-partonic-dr}，第~8~章）。
\uses{lem:prelim-scaleless,def:prelim-dr-conventions}
\end{example}

在计算部分子系数函数时，正是 (ii) 中的红外极点与部分子截面的共线极点相消；弄错它们的来源是共线减除中符号错误的经典来源（本节末的常见误解）。维数正规化并不"消除"红外发散，只是把它也表示成 $1/\varepsilon$ 极点。

\begin{definition}[$\gamma_5$ 与 $\epsilon$ 张量]\label{def:prelim-gamma5}
$\gamma_5=\ii\gamma^0\gamma^1\gamma^2\gamma^3$，Levi-Civita 张量取 $\epsilon^{0123}=+1$（于是 $\epsilon_{0123}=-1$），从而
\begin{equation}\label{eq:prelim-trace-g5}
\Tr(\gamma^\mu\gamma^\nu\gamma^\rho\gamma^\sigma\gamma_5)=-4\ii\,\epsilon^{\mu\nu\rho\sigma}. % VF-CONV-10
\end{equation}
横向张量为 $g_\perp^{\mu\nu}=g^{\mu\nu}-n_A^\mu n_B^\nu-n_B^\mu n_A^\nu$ 与 $\epsilon_\perp^{\mu\nu}=\epsilon^{\mu\nu\rho\sigma}n_{A\rho}n_{B\sigma}$，$\epsilon_\perp^{12}=+1$。在维数正规化中采用 HVBM 方案（'t Hooft--Veltman~\cite{tHooftVeltman1972}、Breitenlohner--Maison~\cite{BreitenlohnerMaison1977a}）并配合 Larin~\cite{Larin:1993tq}对轴矢流的有限重整化（附录~I）；在本书约定下，Larin 方案中轴矢流的写法为
\begin{equation}\label{eq:prelim-larin}
\gamma^\mu\gamma_5=-\frac{\ii}{6}\,\epsilon^{\mu\nu\rho\sigma}\gamma_\nu\gamma_\rho\gamma_\sigma . % VF-CONV-24
\end{equation}
\uses{def:qcd-dimreg}
\end{definition}

\cref{eq:prelim-trace-g5,eq:prelim-larin}在 $4$ 维中已对全部指标组合用显式 $4\times4$ 矩阵核对。文献中常见的 $\gamma_\mu\gamma_5=\frac{\ii}{6}\epsilon_{\mu\nu\rho\sigma}\gamma^\nu\gamma^\rho\gamma^\sigma$ 以 $\epsilon_{0123}=+1$ 为前提，与本书相反，照抄会得到错误的整体符号；引用 $\epsilon$ 张量公式时都要先核对约定（例如 Vanderhaeghen、Guichon 与 Guidal~\cite{Vanderhaeghen:1999xj}采用 $\epsilon_{0123}=+1$）。
% -------------------------------------------------------------------------------------
\section{$[0,1]$ 上的分布与 Mellin 变换}\label{sec:prelim-distributions}

单圈实辐射的相空间中处处出现 $(1-z)^{-1-2\varepsilon}$（$z$ 为纵向动量分数，$z\to1$ 时辐射的胶子变软）。$\varepsilon\to0$ 时它在 $z=1$ 处不可积，只有与虚修正的 $\delta(1-z)$ 项合在一起才有意义；分裂函数与系数函数也在 $z=1$ 处奇异。正确的语言是 $[0,1]$ 上的分布，即作用在检验函数（部分子分布）上的线性泛函。

\begin{definition}[加号分布]\label{def:prelim-plus}
(a) 设 $g$ 在每个 $[0,1-\delta]$ 上可积，且在 $z\to1$ 处至多有 $|\ln^k(1-z)|/(1-z)$ 型的奇异性。对在 $[0,1]$ 上连续、在 $z=1$ 处 Hölder 连续（存在 $\eta>0$ 使 $|h(z)-h(1)|\le C(1-z)^\eta$）的检验函数 $h$，定义
\begin{equation}\label{eq:prelim-plus}
\int_0^1\dd z\,[g(z)]_+\,h(z)=\int_0^1\dd z\,g(z)\big[h(z)-h(1)\big],\qquad \mathcal D_k(z)\equiv\Big[\frac{\ln^k(1-z)}{1-z}\Big]_+ . % VF-CONV-7
\end{equation}
特别地 $\int_0^1[g]_+=0$；$k=0$ 时也写 $\mathcal D_0=[1/(1-z)]_+=1/(1-z)_+$。
\par
(b) 乘积约定：对在 $z=1$ 处 Hölder 连续的函数 $g$，
\begin{equation}\label{eq:prelim-product}
g(z)\,\mathcal D_k\equiv g(1)\,\mathcal D_k+\frac{g(z)-g(1)}{1-z}\ln^k(1-z),
\end{equation}
即 $\langle g\mathcal D_k,h\rangle\equiv\langle\mathcal D_k,gh\rangle$；$k=0$ 时写作 $g(z)/(1-z)_+$。
\par
(c) 常用恒等式：
\begin{gather}
C_F\Big[\frac{1+z^2}{1-z}\Big]_+=C_F\Big[\frac{1+z^2}{(1-z)_+}+\frac32\delta(1-z)\Big],\qquad
\Big[\frac{2z}{1-z}\Big]_+=\frac{2z}{(1-z)_+}+2\delta(1-z),\notag\\
\Big[\frac{z}{1-z}\Big]_+=z\Big[\frac1{1-z}\Big]_++\delta(1-z),\qquad
(1+z^2)\,\mathcal D_1=\Big[\frac{(1+z^2)\ln(1-z)}{1-z}\Big]_++\frac74\,\delta(1-z). \label{eq:prelim-plus-identities} % VF-CONV-7, VF-DY-12
\end{gather}
% uses: none
\end{definition}

(c) 中的四个恒等式都是下面一般公式的特例：对在 $z=1$ 处 Hölder 连续的 $g$，
\begin{equation}\label{eq:prelim-ck}
g(z)\,\mathcal D_k=\Big[\frac{g(z)\ln^k(1-z)}{1-z}\Big]_++c_k[g]\,\delta(1-z),\qquad c_k[g]=\int_0^1\dd z\,\frac{g(z)-g(1)}{1-z}\ln^k(1-z),
\end{equation}
（两边作用于同一个 $h$ 后相减，差为 $h(1)c_k[g]$。）于是 $c_0[1+z^2]=-\int_0^1(1+z)=-\frac32$，$c_0[2z]=-2$，$c_0[z]=-1$，$c_1[1+z^2]=-\int_0^1(1+z)\ln(1-z)=\frac74$。最后一式是实际的陷阱：Drell--Yan 单圈系数函数中的 $4C_F(1+z^2)\mathcal D_1$ 若误作整体加号分布，$\delta(1-z)$ 的系数会差 $7C_F$。检验函数类与乘积约定的完整讨论见\cref{sec:appA-plus}；事件形状中在 $\tau=0$ 处减除的加号分布与 SCET 的星分布有各自的记号（第~15~章），从不与\cref{eq:prelim-plus}混用。

\begin{lemma}[$(1-z)^{-1-a\varepsilon}$ 的分布展开]\label{lem:prelim-plus-expansion}
设 $a>0$，检验函数 $h$ 在 $[0,1]$ 上连续，且存在 $\eta>0$ 与 $C$ 使 $|h(z)-h(1)|\le C(1-z)^\eta$。则对 $-\eta<a\varepsilon<0$，
\begin{equation}\label{eq:prelim-plus-expansion}
\int_0^1\dd z\,(1-z)^{-1-a\varepsilon}h(z)=-\frac{h(1)}{a\varepsilon}+\sum_{k=0}^\infty\frac{(-a\varepsilon)^k}{k!}\int_0^1\dd z\,\mathcal D_k(z)h(z),
\end{equation}
级数绝对收敛，截断到 $k=K$ 的余项为 $O(\varepsilon^{K+1})$。即作为分布
\[
(1-z)^{-1-a\varepsilon}=-\frac1{a\varepsilon}\,\delta(1-z)+\sum_{k\ge0}\frac{(-a\varepsilon)^k}{k!}\,\mathcal D_k(z).
\]
DIS 中用 $a=1$，Drell--Yan 与软函数中用 $a=2$。
\uses{def:prelim-plus}
\end{lemma}

\begin{proof}
完整证明见\cref{sec:appA-plus}。要点：令 $t=-a\varepsilon\in(0,\eta)$，$h=h(1)+[h-h(1)]$；第一部分给出 $h(1)/t$，第二部分代入 $(1-z)^{t}=\sum_kt^k\ln^k(1-z)/k!$ 逐项积分，控制函数 $C(1-z)^{\eta-|t|-1}$ 可积。
\end{proof}

右端对 $|a\varepsilon|<\eta$ 定义了只在 $\varepsilon=0$ 处有单极点的亚纯函数，即左端从红外收敛的一侧 $a\varepsilon<0$ 出发的解析延拓；软胶子的发散正以这个 $-\frac1{a\varepsilon}\delta(1-z)$ 极点出现。

\begin{example}[单圈实辐射的典型积分]\label{ex:prelim-plus}
考虑 $I(\varepsilon)=\int_0^1\dd z\,(1-z)^{-1-2\varepsilon}z^{-\varepsilon}h(z)$。对 $H(z)\equiv z^{-\varepsilon}h(z)$（$H(1)=h(1)$）用\cref{lem:prelim-plus-expansion}（$a=2$），再展开 $H=h-\varepsilon h\ln z+O(\varepsilon^2)$，并注意 $\langle\mathcal D_0,h\ln z\rangle=\int_0^1h\ln z/(1-z)$（因为 $\ln1=0$），得
\begin{equation}\label{eq:prelim-ex-plus}
I(\varepsilon)=-\frac{h(1)}{2\varepsilon}+\langle\mathcal D_0,h\rangle-\varepsilon\Big[2\langle\mathcal D_1,h\rangle+\int_0^1\dd z\,\frac{h(z)\ln z}{1-z}\Big]+O(\varepsilon^2),
\end{equation}
即作为分布 $(1-z)^{-1-2\varepsilon}z^{-\varepsilon}=-\frac{\delta(1-z)}{2\varepsilon}+\mathcal D_0-\varepsilon\big[2\mathcal D_1+\frac{\ln z}{1-z}\big]+O(\varepsilon^2)$。

核对：取 $h(z)=z$，左端有闭式 $\int_0^1(1-z)^{-1-2\varepsilon}z^{1-\varepsilon}\dd z=B(2-\varepsilon,-2\varepsilon)=\Gamma(2-\varepsilon)\Gamma(-2\varepsilon)/\Gamma(2-3\varepsilon)$。右端的三个系数为 $\langle\mathcal D_0,z\rangle=\int_0^1\frac{z-1}{1-z}\dd z=-1$，$\langle\mathcal D_1,z\rangle=-\int_0^1\ln(1-z)\,\dd z=1$，$\int_0^1\frac{z\ln z}{1-z}\dd z=-\sum_{j\ge0}(2+j)^{-2}=1-\zeta_2$，于是\cref{eq:prelim-ex-plus}给出 $-\frac1{2\varepsilon}-1-(3-\zeta_2)\varepsilon=-\frac1{2\varepsilon}-1-1.3551\,\varepsilon$，与 $\Gamma$ 函数闭式的直接展开一致。这是一切单圈实辐射的结构：软极点 $-\frac{1}{2\varepsilon}\delta(1-z)$ 与虚修正的 $\delta(1-z)$ 极点相消，留下 $\mathcal D_k$ 与正规函数（Drell--Yan 的实辐射带 $z^{+\varepsilon}$，只改变 $\ln z$ 项的符号）。
\uses{lem:prelim-plus-expansion}
\end{example}

演化方程是 $x$ 空间的卷积，在 Mellin 空间中化为乘法；阈值区 $z\to1$ 对应大 $N$。

\begin{definition}[Mellin 变换与卷积]\label{def:prelim-mellin}
(a) $[0,1]$ 上函数（或分布）$f$ 的 Mellin 变换为 $\tilde f(N)=\int_0^1\dd z\,z^{N-1}f(z)$（对分布即 $\langle f,z^{N-1}\rangle$；例如 $\delta(1-z)$ 的 Mellin 变换为 $1$）。
\par
(b) 卷积 $(f\otimes g)(x)=\int_x^1\frac{\dd\xi}{\xi}\,f(\xi)\,g(x/\xi)=\int_0^1\!\dd\xi\int_0^1\!\dd y\,f(\xi)g(y)\,\delta(x-\xi y)$。
\par
(c) 调和和 $S_k(N)=\sum_{j=1}^Nj^{-k}$，其解析延拓 $S_1(N)=\psi(N+1)+\gamma_E$，$S_2(N)=\zeta_2-\psi'(N+1)$；$\bar N=N\ee^{\gamma_E}$。% VF-CONV-8, VF-MISC-1
\uses{def:prelim-plus}
\end{definition}

\begin{lemma}[Mellin 卷积定理与加号分布的矩]\label{lem:prelim-mellin}
(a) 对 $[0,1]$ 上的可积函数，以及它们与 $\delta(1-z)$、$g\,\mathcal D_k$ 的有限线性组合（卷积按 $\langle f\otimes g,h\rangle\equiv\langle f(\xi),\langle g(y),h(\xi y)\rangle\rangle$ 定义），$\widetilde{f\otimes g}(N)=\tilde f(N)\,\tilde g(N)$。
\par
(b) 对整数 $N\ge1$，
\begin{equation}\label{eq:prelim-mellin-D}
\int_0^1\dd z\,z^{N-1}\mathcal D_0(z)=-S_1(N-1),\qquad \int_0^1\dd z\,z^{N-1}\mathcal D_1(z)=\frac12\big[S_1^2(N-1)+S_2(N-1)\big], % VF-CONV-8
\end{equation}
并经\cref{def:prelim-mellin}(c) 延拓到 $\Real N>0$。$N\to\infty$ 时二者分别趋于 $-\ln\bar N+O(1/N)$ 与 $\frac12(\ln^2\bar N+\zeta_2)+O(\ln N/N)$。
\uses{def:prelim-mellin,def:prelim-plus,lem:appA-mellin-table}
\end{lemma}

\begin{proof}
完整证明见\cref{sec:appA-mellin}（\cref{lem:appA-mellin-table}）。(a) 取 $h(x)=x^{N-1}$，$h(\xi y)=\xi^{N-1}y^{N-1}$ 分解为乘积。(b) 由 $(z^{N-1}-1)/(1-z)=-\sum_{j=0}^{N-2}z^j$ 得第一式；第二式用 $\int_0^1z^j\ln(1-z)\dd z=-S_1(j+1)/(j+1)$ 与 $\sum_{m=1}^{N-1}S_1(m)/m=\frac12[S_1^2+S_2](N-1)$；大 $N$ 极限来自 $\psi(N)=\ln N+O(1/N)$、$\psi'(N)=O(1/N)$。
\end{proof}

用\cref{lem:prelim-mellin}与 $\int_0^1z^{N-1}\frac{\ln z}{1-z}\dd z=-[\zeta_2-S_2(N-1)]$，\cref{ex:prelim-plus}对一般 $N$ 给出 $B(N-\varepsilon,-2\varepsilon)=-\frac1{2\varepsilon}-S_1(N-1)-\varepsilon[S_1^2+2S_2-\zeta_2](N-1)+O(\varepsilon^2)$，与 $\Gamma$ 函数闭式逐阶一致。大 $N$ 近似收敛较慢：$N=100$ 时 $-S_1(99)=-5.1774$，$-\ln\bar N=-5.1824$；$\frac12[S_1^2+S_2](99)=14.2201$，$\frac12(\ln^2\bar N+\zeta_2)=14.2510$。阈值重求和（第~13~章）所用的正是这些对数。

% -------------------------------------------------------------------------------------
\section{矩问题、Carlson 定理与 Müntz--Szász 定理}\label{sec:prelim-moments}

严格核心中的深度非弹散射定理（第~7~章）给出的是结构函数的矩，演化在 Mellin 空间中是对角的。于是要问：矩能否唯一地确定 $x$ 空间的函数？整数（甚至只是偶数）$N$ 上的值能否唯一地延拓到复 $N$ 平面？附录~K 证明了相应的引理；这里汇总结论，并说明它们\emph{不能}回答什么。

\begin{remark}[矩问题工具箱]\label{rem:prelim-moment-tools}
\begin{enumerate}[label=(\roman*),leftmargin=2.2em,itemsep=0pt,topsep=2pt]
\item \emph{Hausdorff 唯一性}：$[0,1]$ 上全部矩为零的有限带号 Borel 测度为零；全部矩为零的 $L^1$ 函数几乎处处为零（\cref{lem:appK-hausdorff}）。
\item \emph{偶数矩已足够}：若 $x^{2k_0}g\in L^1(0,1)$ 且 $\int_0^1x^{2k}g\,\dd x=0$ 对一切 $k\ge k_0$ 成立，则 $g=0$ 几乎处处成立（\cref{thm:appK-even-moments}）。对部分子分布，偶数矩（$N\ge N_1$）确定 $q+\bar q$ 与 $g$，奇数矩确定 $q-\bar q$，单靠偶数矩不能把 $q$ 与 $\bar q$ 分开（\cref{cor:appK-pdf}）。
\item \emph{延拓的唯一性}：Carlson 定理（\cref{thm:appK-carlson}）——在 $\Real z\ge0$ 上全纯、指数型、沿虚轴增长慢于 $\ee^{c|y|}$（$c<\pi$）且在全部非负整数上为零的函数恒为零（$\pi$ 不可改进，反例 $\sin\pi z$；只用偶数点时界为 $\pi/2$，反例 $\sin(\pi z/2)$）。对矩而言更简单的版本已足够：在 $\Real z>\sigma_0$ 上全纯、多项式有界且在一个算术级数上为零的函数恒为零（\cref{prop:appK-continuation}）。所以 Mellin 变换由偶数 $N$ 唯一地延拓（例如 $S_1(N)\mapsto\psi(N+1)+\gamma_E$；加上 $c\sin(\pi N/2)$ 违反增长条件），而由偶数与由奇数 $N$ 出发的延拓一般不同（分别对应 $q+\bar q$ 与 $q-\bar q$）。
\item \emph{Müntz--Szász 定理}：对 $0<\lambda_1<\lambda_2<\dots\to\infty$，$\{1,x^{\lambda_1},x^{\lambda_2},\dots\}$ 的线性组合在 $C[0,1]$ 中稠密，当且仅当 $\sum_k1/\lambda_k=\infty$（\cref{thm:appK-muntz}）。偶数幂满足此条件，但 (ii) 并不需要它。
\item \emph{反例与不稳定性}（\cref{rem:appK-counterexamples}）：$\delta'(x)$ 的全部偶数矩为零，却不是零分布（它不是测度）；$R_Q(x)=\sin(Q^2x)$ 对每个固定的 $n$ 满足 $|Q^2\int_0^1x^nR_Q\,\dd x|\le c_n$，但 $\sup|R_Q|=1$；唯一性不等于稳定性——矩的 Gram（Hilbert）矩阵的条件数 $\approx\ee^{3.5K}$，$K=5$ 时为 $4.8\times10^5$，$K=20$ 时为 $2.5\times10^{28}$。% VF-MISC-2, VF-MISC-3
\end{enumerate}
结论：矩只能把 $x$ 空间的对象确定为分布（测度）；由此既得不到逐点的界，也得不到稳定的反演。
\uses{lem:appK-hausdorff,thm:appK-even-moments,cor:appK-pdf,thm:appK-carlson,prop:appK-continuation,thm:appK-muntz,rem:appK-counterexamples}
\end{remark}

由 (i)、(ii) 与附录~A 的弱收敛引理，算符乘积展开对矩的 $O(1/Q^2)$ 控制只推出 $x$ 空间的弱收敛（\cref{cor:dis-weak-xspace}，第~7~章）；(v) 说明逐点的 $O(1/Q^2)$ 无从谈起（\cref{rem:dis-nonuniform}），逐点陈述只能来自图论的全阶论证（第~10~章）。(iii) 中偶、奇延拓的区别是非单态演化核 $P^\pm_{\mathrm{NS}}$ 从次领头阶起不同的数学根源（第~8~章）；(v) 中的条件数说明由有限个格点矩重构 $x$ 依赖为什么是病态问题（第~22~章）。
% -------------------------------------------------------------------------------------
\section{割图、切割规则、幺正性与光学定理}\label{sec:prelim-cuts}

截面是振幅模平方对末态的求和，而因子化证明处理的是前向振幅的 Feynman 图；两者由幺正性联系：一个图的全部割之和等于它的虚部。规范理论中还有第二个问题：割胶子线用 $-g_{\mu\nu}$ 求极化和会把非物理极化也算进去，需要鬼场补偿。本节给出割图、切割规则与物理幺正性，证明流关联函数的光学定理，并用单圈真空极化把它们串起来。

\begin{definition}[割图]\label{def:prelim-cut}
一个割图是一个 Feynman 图被一条"割线"分成左、右两部分，割线把入射外线与出射外线分开，并配以如下规则：
\begin{enumerate}[label=(\alph*),leftmargin=2.2em,itemsep=0pt,topsep=2pt]
\item 割线左侧（振幅）用通常的 Feynman 规则（\cref{def:prelim-feynman-rules}），传播子带 $+\ii0$；
\item 割线右侧（共轭振幅）用复共轭的规则：传播子取复共轭（$-\ii0$），顶点取复共轭（例如夸克--胶子顶点 $\ii g\gamma^\mu t^a\to-\ii g\gamma^\mu t^a$；Wilson 线的规则见\cref{prop:prelim-wilson-rules}）；
\item 被割的线换成在壳、正能量的割线因子：末态夸克 $2\pi\delta^+(p^2-m^2)(\sla p+m)$，末态反夸克 $2\pi\delta^+(p^2-m^2)(\sla p-m)$（$p$ 为该粒子的物理动量，$\delta^+(p^2-m^2)\equiv\theta(p^0)\delta(p^2-m^2)$），胶子 $2\pi\delta^+(k^2)(-g^{\mu\nu})\delta^{ab}$（Feynman 规范），鬼场 $2\pi\delta^+(k^2)\delta^{ab}$，每个闭合鬼圈另带因子 $-1$；
\item 对割线上的动量积分，对末态的颜色、自旋求和，乘以对称因子。
\end{enumerate}
\uses{def:prelim-feynman-rules,prop:prelim-wilson-rules}
\end{definition}

(c) 中的因子来自 $\sum_su\bar u=\sla p+m$、$\sum_sv\bar v=\sla p-m$ 与 $\int\frac{\dd^3p}{(2\pi)^32E}=\int\frac{\dd^4p}{(2\pi)^4}2\pi\delta^+(p^2-m^2)$；按这些规则，单举截面 $\sum_X|\langle X|T|i\rangle|^2$ 就是全部割图之和。部分子分布（第~8~章）也是割图：割线两侧各有一个算符顶点，左侧的 Wilson 线用振幅规则，右侧用其复共轭。

\begin{theorem}[切割规则 \tagR]\label{thm:prelim-cutting-rules}
（$R_{\mathrm{PT}}$：在重整化微扰论中对红外正规化的部分子态严格。）设 $G$ 是一个 Feynman 图，全部内线带正质量（或以其他方式正规化，使所涉积分都收敛），外动量为实数且位于物理区域；把图 $G$ 的值记作 $\ii\mathcal M_G$。则
\begin{equation}\label{eq:prelim-cutting}
2\,\Imag\mathcal M_G=\sum_{c}\mathcal M_{G,c},
\end{equation}
求和遍及把入射外线与出射外线分开、且与外动量的能量流相容的全部割 $c$（Cutkosky 割），$\mathcal M_{G,c}$ 是按\cref{def:prelim-cut}计算的割图。在坐标空间中等价的陈述是最大时间方程：对图的全部 $2^n$ 种顶点标记求和为零。（本定理只涉及 Feynman 图。）
\uses{def:prelim-cut,def:prelim-feynman-rules}
\end{theorem}

\begin{proofsketch}
(1) 坐标空间：把每个顶点标记为"无圈"或"有圈"（有圈顶点用复共轭的顶点因子），两端无圈的线用 $\Delta_F(x-y)$，两端有圈的用 $\Delta_F^*(x-y)$，无圈的 $x$ 与有圈的 $y$ 之间用正频函数 $\Delta^+(y-x)$。由 $\Delta_F(x)=\theta(x^0)\Delta^+(x)+\theta(-x^0)\Delta^+(-x)$，若 $x_i$ 是时间最大的顶点，与它相连的线在 $x_i$ 改变标记时都不变而顶点因子变号；故全部 $2^n$ 种标记之和为零。

(2) 全无圈的项是 $\ii\mathcal M_G$，全有圈的项是其复共轭 $-\ii\mathcal M_G^*$，二者之和为 $-2\Imag\mathcal M_G$；由 (1)，$2\Imag\mathcal M_G$ 等于全部混合标记之和。

(3) 动量空间中正频函数是 $2\pi\theta(k^0)\delta(k^2-m^2)$（乘以分子），正能量只能由无圈区流向有圈区；对给定的外动量，只有无圈区含全部入射外线、有圈区含全部出射外线的标记保留下来，这正是 Cutkosky 割，有圈区的规则即\cref{def:prelim-cut}的右侧。完整证明见 Veltman~\cite{Veltman1963}、'tHooft--Veltman~\cite{tHooftVeltman1974}与 Cutkosky~\cite{Cutkosky1960}；教科书式的讲解见\cite[第~7.3~节]{PeskinSchroeder1995}。
\end{proofsketch}

附录~O 用到这一方程的光锥版本：带算符插入时移动割线位置的恒等式及其闭时路形式（\cref{lem:appO-cut-position}、\cref{lem:appO-unitarity}）。Feynman 规范中割胶子线的 $-g^{\mu\nu}$ 含非物理极化，下一命题说明鬼场割线恰好把它们抵消。

\begin{proposition}[物理幺正性与鬼场割线 \tagR]\label{prop:prelim-unitarity}
（$R_{\mathrm{PT}}$：在重整化微扰论中对红外正规化的部分子态严格。）在 Feynman 规范、维数正规化的微扰 QCD 中，考虑一个前向振幅任一阶的割图，并设所涉割图积分有定义（红外已正规化）。对每一组固定的割线动量，割胶子线以 $-g_{\mu\nu}$ 求和并计入全部鬼场割线（每个闭合鬼圈带 $-1$）所得之和，等于只对物理末态（夸克、反夸克与横向极化胶子）求和所得之和。结合\cref{thm:prelim-cutting-rules}，前向振幅的虚部逐阶等于对物理末态的截面之和。树图例子 $q\bar q\to gg$（\cref{ex:qcd-qqbar-gg}）中，% VF-WL-9
\begin{equation}\label{eq:prelim-qqgg}
\sum_{\mathrm{phys}}|\mathcal A|^2=\sum_{-g,-g}|\mathcal A|^2-2\sum|\mathcal A_{\mathrm{gh}}|^2,\qquad
\frac1{36}\sum_{\mathrm{phys}}|\mathcal A|^2=g^4\Big[\frac{32}{27}\frac{t^2+u^2}{tu}-\frac83\frac{t^2+u^2}{s^2}\Big],
\end{equation}
因子 $2$ 来自鬼与反鬼两种末态，$\frac1{36}$ 是夸克自旋与颜色的平均。（本命题只涉及部分子态。）
\uses{def:prelim-cut,thm:prelim-cutting-rules,prop:qcd-brst-cohomology,thm:prelim-slavnov-taylor,prop:appJ-ghost-cuts,ex:qcd-qqbar-gg}
\end{proposition}

\begin{proofsketch}
(1) 对在壳胶子取类光的 $\bar k=(k^0,-\mathbf k)$，$-g_{\mu\nu}=\sum_{\lambda=1,2}\epsilon^{(\lambda)}_\mu\epsilon^{(\lambda)*}_\nu-\frac{k_\mu\bar k_\nu+\bar k_\mu k_\nu}{k\cdot\bar k}$。(2) 由 Slavnov--Taylor 恒等式（\cref{thm:prelim-slavnov-taylor}），把割胶子的极化换成其动量时振幅剩下鬼场振幅（树图形式 $k_{1\mu}\mathcal A^{\mu\nu}_{ab}=\frac{g^2f^{abc}}{Q^2}(k_2^2g^{\nu\sigma}-k_2^\nu k_2^\sigma)\bar v(p_2)\gamma_\sigma t^cu(p_1)$，另一胶子为物理极化时为零），于是 (1) 中第二项与鬼场割线相消。(3) 一般的论证是 BRST 四重态机制（\cref{prop:qcd-brst-cohomology}）：非物理模式组成零模四重态，在割线求和中整体抵消。完整证明见\cref{prop:appJ-ghost-cuts}，以及 Kugo--Ojima~\cite{KugoOjima1979}、'tHooft--Veltman~\cite{tHooftVeltman1972b}；树图截面与 Combridge、Kripfganz、Ranft~\cite{CombridgeKripfganzRanft1977}一致。
\end{proofsketch}

\begin{remark}[无质量维数正规化中的割图规则 \tagA]\label{rem:prelim-cutting-dr}
对无质量 QCD，\cref{thm:prelim-cutting-rules}与\cref{prop:prelim-unitarity}推广为维数正规化积分之间的恒等式：对含鬼场的无质量图，作为 $\varepsilon$ 的亚纯函数，对割线求和等于前向振幅的虚部，并等于对物理末态的求和。
% errtype: E-S
\uses{thm:prelim-cutting-rules,prop:prelim-unitarity,def:qcd-dimreg}
\end{remark}

\paragraph{地位与缺口}
在 $O(g^2)$ 处，对流关联函数，这一恒等式已精确到 $\varepsilon$ 的一切阶验证：
\[
2(D-1)\,s\,\Imag\Pi^{(1)}_{\mathrm{bare}}(s+\ii0)=\rho_V+\rho_R , % VF-EE-8
\]
且逐项成立：双圈图中 $G(1,1)^2$ 一项对应虚修正（两部分子割），$G(1,1)G(1+\varepsilon,1)$ 一项对应实辐射（三部分子割）（\cref{prop:ee-cutting-exact}，第~6~章）。
一般情形的论证是标准的：最大时间方程在 $D$ 维被积函数的层次上形式不变，鬼场与非物理极化的抵消是 BRST 四重态的代数性质（Veltman~\cite{Veltman1963}，'tHooft--Veltman~\cite{tHooftVeltman1972b,tHooftVeltman1974}，Kugo--Ojima~\cite{KugoOjima1979}）。最弱的一步是：单个无质量割图有红外发散，逐图应用最大时间方程、再把每个割图作为 $\varepsilon$ 的亚纯函数延拓，需要对这些积分的存在性与延拓作额外论证，本书没有给出。因子化真正需要的精确形式是附录~N 的割线求和公共表示（\cref{asm:appN-cutsum}，\tagA，此处只作指引）。误差类型：已证明之处恒等式是精确的，一般陈述只是 E-S 型的论证。要提升为 [R]，需要对无质量、含鬼场的维数正规化图证明逐图的最大时间方程及其解析延拓。

\begin{proposition}[流关联函数的光学定理 \tagR]\label{prop:prelim-optical}
（$R_{\mathrm{def}}$：由算符定义直接推出。）设 $j^\mu=\sum_qe_q\bar q\gamma^\mu q$（以 $e$ 为单位），对自旋平均的靶态 $|P\rangle$ 令 % VF-CONV-12, VF-CONV-13, VF-CONV-14
\begin{gather*}
\Pi^{\mu\nu}(q)=\ii\int\dd^4x\,\ee^{\ii q\cdot x}\langle0|T\,j^\mu(x)j^\nu(0)|0\rangle=(q^\mu q^\nu-q^2g^{\mu\nu})\,\Pi(q^2),\\
W^{\mu\nu}=\frac1{4\pi}\int\dd^4z\,\ee^{\ii q\cdot z}\langle P|[j^\mu(z),j^\nu(0)]|P\rangle,\qquad
T^{\mu\nu}=\ii\int\dd^4z\,\ee^{\ii q\cdot z}\langle P|T\,j^\mu(z)j^\nu(0)|P\rangle ,
\end{gather*}
并以 $\Imag X^{\mu\nu}\equiv\frac1{2\ii}[X^{\mu\nu}-(X^{\nu\mu})^*]$ 表示吸收部分（对实对称的张量结构即通常的虚部）。
\par
(a) 对 $q^0>0$，
\begin{equation}\label{eq:prelim-optical-vac}
\sum_X(2\pi)^4\delta^4(q-p_X)\langle0|j^\mu(0)|X\rangle\langle X|j^\nu(0)|0\rangle=2\,\Imag\Pi^{\mu\nu}(q),
\end{equation}
因而 $\Imag\Pi(s+\ii0)\ge0$；$e^+e^-\to$ 强子（单光子交换、无质量轻子、$\alpha$ 的最低阶）的截面为 $\sigma=\frac{16\pi^2\alpha^2}{s}\Imag\Pi(s+\ii0)$，从而 $R(s)\equiv\sigma/\sigma_{\mu\mu}=12\pi\,\Imag\Pi(s+\ii0)$，$\sigma_{\mu\mu}=\frac{4\pi\alpha^2}{3s}$。
\par
(b) 对 $P\cdot q>0$（靶静止系中 $q^0>0$），$W^{\mu\nu}=\frac1{2\pi}\Imag T^{\mu\nu}$，即结构函数满足 $\Imag T_i=2\pi F_i$；$x>1$ 时两边都为零。
\par
非微扰地位：\cref{asm:intro-wightman}（MA0：态的完备性与谱条件）——只对强子态需要；对部分子态，(a)(b) 是微扰论逐阶的恒等式。
\uses{def:prelim-cut,asm:intro-wightman}
\end{proposition}

\begin{proof}
(a) 由平移不变性 $j(x)=\ee^{\ii P\cdot x}j(0)\ee^{-\ii P\cdot x}$ 与 $\langle0|P^\mu=0$，在 $T$ 乘积的两个时序中插入完备态：
\[
\langle0|j^\mu(x)j^\nu(0)|0\rangle=\sum_XA_X^{\mu\nu}\ee^{-\ii p_X\cdot x},\qquad \langle0|j^\nu(0)j^\mu(x)|0\rangle=\sum_XB_X^{\nu\mu}\ee^{+\ii p_X\cdot x},
\]
$A_X^{\mu\nu}=\langle0|j^\mu(0)|X\rangle\langle X|j^\nu(0)|0\rangle$，$B_X^{\nu\mu}=\langle0|j^\nu(0)|X\rangle\langle X|j^\mu(0)|0\rangle$。用
\begin{gather*}
\int\dd^4x\,\ee^{\ii q\cdot x}\theta(x^0)\ee^{-\ii p_X\cdot x}=\frac{\ii\,(2\pi)^3\delta^3(\mathbf q-\mathbf p_X)}{q^0-p_X^0+\ii0},\\
\int\dd^4x\,\ee^{\ii q\cdot x}\theta(-x^0)\ee^{+\ii p_X\cdot x}=\frac{-\ii\,(2\pi)^3\delta^3(\mathbf q+\mathbf p_X)}{q^0+p_X^0-\ii0},
\end{gather*}
得
\[
\Pi^{\mu\nu}(q)=-\sum_X\frac{(2\pi)^3\delta^3(\mathbf q-\mathbf p_X)}{q^0-p_X^0+\ii0}A_X^{\mu\nu}+\sum_X\frac{(2\pi)^3\delta^3(\mathbf q+\mathbf p_X)}{q^0+p_X^0-\ii0}B_X^{\nu\mu}.
\]
流是厄米的，故 $(A_X^{\nu\mu})^*=A_X^{\mu\nu}$，$(B_X^{\mu\nu})^*=B_X^{\nu\mu}$。由 $\frac1{\Delta-\ii0}-\frac1{\Delta+\ii0}=2\pi\ii\,\delta(\Delta)$，第一项的吸收部分为 $\frac12\sum_X(2\pi)^4\delta^4(q-p_X)A_X^{\mu\nu}$；第二项的吸收部分正比于 $\delta(q^0+p_X^0)$，由谱条件 $p_X^0\ge0$ 在 $q^0>0$ 时为零；$T$ 乘积在 $x=0$ 处的定义只差局域项，给出 $q$ 的厄米多项式，没有吸收部分。这就是\cref{eq:prelim-optical-vac}。在 $q$ 的静止系中取空间指标 $\mu=\nu=i$，左边 $\ge0$，右边为 $2s\Imag\Pi$，故 $\Imag\Pi\ge0$。截面：自旋平均的轻子张量 $L_{\mu\nu}=k_{1\mu}k_{2\nu}+k_{1\nu}k_{2\mu}-g_{\mu\nu}k_1\cdot k_2$ 满足 $L_{\mu\nu}(q^\mu q^\nu-q^2g^{\mu\nu})=s^2$，故 $\sigma=\frac1{2s}\frac{e^4}{s^2}L_{\mu\nu}2\Imag\Pi^{\mu\nu}=\frac{16\pi^2\alpha^2}{s}\Imag\Pi$，除以 $\sigma_{\mu\mu}$ 即得 $R$。

(b) 同样插入完备态。对易子的第二项与 $T$ 乘积的 $\theta(-z^0)$ 项要求 $p_X=P-q$，而 $(P-q)^2=M^2-2P\cdot q-Q^2<M^2$，靶的超选择扇区中没有这样的态（谱条件；对部分子靶即带同样量子数的态 $p_X^2\ge m^2$），故它们为零。剩下的项给出
\begin{gather*}
W^{\mu\nu}=\frac1{4\pi}\sum_X(2\pi)^4\delta^4(P+q-p_X)\langle P|j^\mu|X\rangle\langle X|j^\nu|P\rangle,\\
\Imag T^{\mu\nu}=\frac12\sum_X(2\pi)^4\delta^4(P+q-p_X)\langle P|j^\mu|X\rangle\langle X|j^\nu|P\rangle,
\end{gather*}
后者的计算与 (a) 相同；两式相比即得 $W^{\mu\nu}=\frac1{2\pi}\Imag T^{\mu\nu}$，按同一组张量结构分解即 $\Imag T_i=2\pi F_i$。$x>1$ 时 $(P+q)^2=M^2+2P\cdot q(1-x)<M^2$，两边都没有中间态。
\end{proof}

检验：自由夸克圈给出 $\Imag\Pi=N_ce_q^2/(12\pi)$，从而 $R_0=N_c\sum_qe_q^2$（$n_f=5$ 时为 $11/3$），见下例。% VF-EE-3

\begin{example}[维数正规化中的单圈无质量流关联函数]\label{ex:prelim-vacpol}
单个味道（电荷 $e_q$）的单圈图带闭合夸克圈的因子 $-1$ 与色因子 $N_c$，
\[
\Pi^{\mu\nu}(q)=\ii N_ce_q^2\int\frac{\dd^Dk}{(2\pi)^D}\,\frac{\Tr\big[\gamma^\mu\sla k\gamma^\nu(\sla k-\sla q)\big]}{(k^2+\ii0)\,((k-q)^2+\ii0)} .
\]
用 $g_{\mu\nu}$ 缩并：左边为 $-(D-1)q^2\Pi$；右边用 $\gamma^\mu\sla a\gamma_\mu=-(D-2)\sla a$ 与 $\Tr[\sla k(\sla k-\sla q)]=2[k^2+(k-q)^2-q^2]$，其中 $k^2$、$(k-q)^2$ 两项各消去一个传播子而成为无标度积分（\cref{lem:prelim-scaleless}）。再用两点积分（\cref{lem:appC-bubble}）$\int\frac{\dd^Dk}{(2\pi)^D}\frac1{(k^2+\ii0)((k-q)^2+\ii0)}=\frac{\ii\,G(1,1)}{(4\pi)^{D/2}}(-q^2-\ii0)^{-\varepsilon}$，得
\begin{equation}\label{eq:prelim-Pi0}
\Pi_0(q^2)=N_ce_q^2\,\frac{2(D-2)}{D-1}\,\frac{G(1,1)}{(4\pi)^{D/2}}\,(-q^2-\ii0)^{-\varepsilon},\qquad G(1,1)=\frac{\Gamma(\varepsilon)\Gamma(1-\varepsilon)^2}{\Gamma(2-2\varepsilon)}. % VF-EE-3
\end{equation}
\par
\emph{$\MSbar$ 重整化。}光子自能中出现的是 $e_0^2\Pi_0=e^2\big[\mu^{2\varepsilon}\ee^{\varepsilon\gamma_E}(4\pi)^{-\varepsilon}\Pi_0\big]$。由 $\ee^{\varepsilon\gamma_E}G(1,1)=\frac1\varepsilon+2+O(\varepsilon)$（$\ee^{\varepsilon\gamma_E}\Gamma(\varepsilon)$ 无常数项，$\Gamma(1-\varepsilon)^2/\Gamma(1-2\varepsilon)=1+O(\varepsilon^2)$，$\Gamma(2-2\varepsilon)=(1-2\varepsilon)\Gamma(1-2\varepsilon)$）与 $\frac{2(D-2)}{D-1}=\frac43\big(1-\frac\varepsilon3+O(\varepsilon^2)\big)$，
\[
\mu^{2\varepsilon}\ee^{\varepsilon\gamma_E}(4\pi)^{-\varepsilon}\Pi_0=\frac{N_ce_q^2}{12\pi^2}\Big[\frac1\varepsilon+\frac53-L+O(\varepsilon)\Big],\qquad L=\ln\frac{-q^2-\ii0}{\mu^2},
\]
减去极点得 $\Pi_0^{\MSbar}=\frac{N_ce_q^2}{12\pi^2}\big(\frac53-L\big)$。
\par
\emph{虚部的两种算法。}(i) $q^2=s+\ii0$ 处 $L=\ln(s/\mu^2)-\ii\pi$，故 $\Imag\Pi_0^{\MSbar}=\frac{N_ce_q^2}{12\pi}$，$R_0=12\pi\Imag\Pi_0=N_ce_q^2$。(ii) 精确到 $\varepsilon$：$(-s-\ii0)^{-\varepsilon}=s^{-\varepsilon}\ee^{\ii\pi\varepsilon}$ 与 $\Gamma(\varepsilon)\Gamma(1-\varepsilon)\sin\pi\varepsilon=\pi$ 给出
\[
2(D-1)\,s\,\Imag\Pi_0(s+\ii0)=N_ce_q^2\,2(D-2)\,s\,\Phi_2,\qquad \Phi_2=\frac1{8\pi}\Big(\frac{4\pi}{s}\Big)^{\varepsilon}\frac{\Gamma(1-\varepsilon)}{\Gamma(2-2\varepsilon)},
\]
右边正是两部分子割：按\cref{def:prelim-cut}，以 $-g_{\mu\nu}$ 缩并的谱函数为 $N_ce_q^2\int\dd\Phi_2(-g_{\mu\nu})\Tr[\sla p_1\gamma^\mu\sla p_2\gamma^\nu]=N_ce_q^2\,2(D-2)s\,\Phi_2$（$\Phi_2$ 为 $D$ 维两体相空间的体积，附录~C）。所以按切割规则（\cref{thm:prelim-cutting-rules}）与光学定理（\cref{prop:prelim-optical}）算出的虚部与直接延拓的虚部对一切 $\varepsilon$ 精确相等。
\par
\emph{符号。}Peskin--Schroeder 的单粒子不可约光子自能写作 $(q^2g^{\mu\nu}-q^\mu q^\nu)\Pi_{\mathrm{PS}}$，与本书的关系为 $\Pi_{\mathrm{PS}}=-e^2\Pi$；本书的符号使 $\Imag\Pi>0$。
\uses{def:prelim-dr-conventions,def:prelim-msbar,thm:prelim-cutting-rules,prop:prelim-optical,lem:appC-bubble}
\end{example}

\begin{tcolorbox}[enhanced,breakable,title=常见误解,fonttitle=\bfseries,
  colback=gray!3!white,colframe=gray!55!black,boxrule=0.6pt,arc=2pt,
  left=5pt,right=5pt,top=3pt,bottom=3pt,before skip=6pt,after skip=6pt]
\textbf{「维数正规化中夸克中的夸克分布就是 $\delta(1-x)$，所以什么也没有发生。」}只对裸分布成立：全部圈图都是无标度的，$\int\dd^{2-2\varepsilon}\mathbf l_{\mathrm T}/\mathbf l_{\mathrm T}^2\propto\frac1{\varepsilon_{\mathrm{UV}}}-\frac1{\varepsilon_{\mathrm{IR}}}=0$。重整化分布是 $\delta(1-x)-\frac{\as}{2\pi}\frac1\varepsilon P_{qq}+\dots$，其中的极点是红外的（\cref{ex:prelim-scaleless}；一般陈述见\cref{prop:pdf-partonic-dr}）。\par
\textbf{「单个割图是幺正的。」}不对。切割规则把前向图的虚部等同于它的\emph{全部}割之和；末态相互作用以及虚修正与实辐射的抵消只在对割线位置求和后发生，单个割图甚至是红外发散的（\cref{prop:ee-cutting-exact}）。因子化中末态相互作用的相消正是这样的割线求和（\cref{prop:method-fsi}，第~9~章；精确形式见\cref{lem:appO-unitarity}）。
\end{tcolorbox}
% -------------------------------------------------------------------------------------
\paragraph{本章小结}
沿 $z$ 轴的推动把光锥分量乘以 $\ee^{\pm\eta}$、把好/坏旋量分量乘以 $\ee^{\pm\eta/2}$（\cref{lem:prelim-boost}），这是幂次计数的出发点。Wilson 线是规范协变的平行移动算符（\cref{prop:prelim-wilson-covariance}），其动量空间规则由一条统一规则给出，未来与过去指向的线相差 $2\pi\ii\,\delta(n\cdot k)$（\cref{prop:prelim-wilson-rules}）；程函恒等式在绝热正规化下精确成立（\cref{lem:prelim-eikonal-identity}），并把高能部分子与软胶子的耦合等同于 Wilson 线（\cref{ex:prelim-eikonal-two}）。维数正规化的每个极点都要标明来源，无标度积分为零（\cref{lem:prelim-scaleless,ex:prelim-scaleless}）。加号分布与 $(1-z)^{-1-a\varepsilon}$ 的展开组织了单圈实辐射（\cref{lem:prelim-plus-expansion}），Mellin 变换把卷积化为乘积（\cref{lem:prelim-mellin}），矩只把 $x$ 空间的对象确定为分布（\cref{rem:prelim-moment-tools}）。切割规则、物理幺正性与光学定理（\cref{thm:prelim-cutting-rules,prop:prelim-unitarity,prop:prelim-optical}）把截面与前向振幅的虚部联系起来；无质量维数正规化中的割图规则是全阶论证（\cref{rem:prelim-cutting-dr}，\tagA）。

\paragraph{这些工具在后文的用途}
光锥坐标与推动引理用于第~4~章的区域标度与幂次计数（\cref{def:ir-scalings}、\cref{prop:ir-power-counting}）。Wilson 线出现在部分子分布的定义（\cref{def:pdf-quark}）、尖点反常维度（\cref{def:sudakov-cusp}）、共线与软近似（\cref{def:method-collinear-approx}、\cref{def:method-soft-approx}）、多胶子 Ward 恒等式（\cref{lem:method-ward-multi}）与 TMD 的定义（\cref{def:tmd-unsub}）中。维数正规化的记账用于\cref{prop:sudakov-oneloop}与\cref{prop:pdf-partonic-dr}，$\gamma_5$ 的约定用于\cref{thm:dis-bjorken}与\cref{thm:coll-g1-moments}。加号分布与 Mellin 变换用于分裂函数（\cref{prop:pdf-lo-kernels}）、次领头阶系数函数（\cref{prop:disall-nlo}、\cref{prop:dy-nlo}）与阈值区（\cref{prop:thr-oneloop}）；矩问题的工具用于\cref{cor:dis-weak-xspace}、\cref{rem:dis-nonuniform}与\cref{rem:pdf-ns-plusminus}。割图与光学定理是\cref{thm:ee-R-finite}、\cref{prop:dis-dispersion}、\cref{prop:method-fsi}与\cref{lem:method-soft-unitarity}的基础，其光锥形式在附录~O 中证明。

\begin{table}[htbp]
\centering
\caption{第~3~章记号速查}\label{tab:prelim-symbols}
\footnotesize
\renewcommand{\arraystretch}{1.05}
\begin{tabularx}{\textwidth}{@{}l X l@{}}
\toprule
记号 & 含义 & 定义处 \\
\midrule
$v^\pm$，$\mathbf v_{\mathrm T}$，$y$ & 光锥分量 $(v^0\pm v^3)/\sqrt2$，横向分量，快度 $\frac12\ln(k^+/k^-)$ & \cref{def:prelim-lightfront} \\
$n_A$，$n_B$ & 类光矢量，$n_A\cdot v=v^-$，$n_B\cdot v=v^+$，$n_A\cdot n_B=1$ & \cref{def:prelim-lightfront} \\
$P_\pm$ & 好/坏分量的投影 $P_+=\frac12\gamma^-\gamma^+$，$P_-=\frac12\gamma^+\gamma^-$ & \cref{lem:prelim-boost} \\
$\lambda$ & 幂次计数参数 $m/Q$；$A$-共线标度 $Q(1,\lambda^2,\lambda)$ & \cref{ex:prelim-boost-scaling} \\
$W[b,a]$ & Wilson 线 $P\exp(\ii g\int_a^b\dd z\cdot A)$，较晚的点在左 & \cref{def:prelim-wilson} \\
$W[\infty n,0]$，$W[0,-\infty n]$ & 未来指向、过去指向的半无限线 & \cref{def:prelim-wilson} \\
$W^{(A)}_{ab}$，$\tilde A(k)$ & 伴随 Wilson 线 $2\Tr[t^aWt^bW^\dagger]$；胶子场的 Fourier 分量（$k$ 流入） & \cref{def:prelim-wilson} \\
$\epsilon$ & 绝热参数（与维数正规化的 $\varepsilon$ 无关） & \cref{def:prelim-wilson} \\
$\varepsilon_{\mathrm{UV}}$，$\varepsilon_{\mathrm{IR}}$ & 维数正规化极点的来源标注 & \cref{def:prelim-dr-conventions} \\
$\gamma_5$，$\epsilon^{\mu\nu\rho\sigma}$ & $\gamma_5=\ii\gamma^0\gamma^1\gamma^2\gamma^3$，$\epsilon^{0123}=+1$ & \cref{def:prelim-gamma5} \\
$g_\perp^{\mu\nu}$，$\epsilon_\perp^{\mu\nu}$ & 横向度规与横向反对称张量，$\epsilon_\perp^{12}=+1$ & \cref{def:prelim-gamma5} \\
$[g]_+$，$\mathcal D_k$ & 加号分布，$\mathcal D_k=[\ln^k(1-z)/(1-z)]_+$ & \cref{def:prelim-plus} \\
$\tilde f(N)$，$\otimes$ & Mellin 变换 $\int_0^1z^{N-1}f$，Mellin 卷积 & \cref{def:prelim-mellin} \\
$S_k(N)$，$\bar N$ & 调和和，$\bar N=N\ee^{\gamma_E}$ & \cref{def:prelim-mellin} \\
$\Pi(q^2)$，$R(s)$ & 流关联函数，$R=12\pi\Imag\Pi(s+\ii0)$ & \cref{prop:prelim-optical} \\
$W^{\mu\nu}$，$T^{\mu\nu}$ & 强子张量与前向 Compton 振幅，$W^{\mu\nu}=\frac1{2\pi}\Imag T^{\mu\nu}$ & \cref{prop:prelim-optical} \\
$\Imag X^{\mu\nu}$ & 吸收部分 $\frac1{2\ii}[X^{\mu\nu}-(X^{\nu\mu})^*]$ & \cref{prop:prelim-optical} \\
\bottomrule
\end{tabularx}
\end{table}
